\documentclass[../main.tex]{subfiles}
\begin{document}
%==========================================TIÊU ĐỀ TẬP========================================
\begin{table}[h]
  \begin{adjustwidth}{-1cm}{}
    \begin{tabular}{>{\centering\arraybackslash}p{6cm}|>{\raggedright\arraybackslash}p{12.5cm}}
      \multirow{5}{6cm}
      % Cột bên trái:  ÉPISODE và số tập
      {\centering
        {\\[-40pt]
          \flaregothic\fontsize{20pt}{20pt}\selectfont
          \color{eptype}ÉPISODE\color{black}\\
          \burbank\fontsize{72pt}{72pt}\selectfont
          \color{epnum}93
          \\[8pt]
          \hrule
          \vspace{8pt}
          \flaregothic\fontsize{20pt}{20pt}\selectfont
          \color{eptype}SPÉCIAL\color{black}\\
          \burbank\fontsize{72pt}{72pt}\selectfont
          \color{epnum}XVII\color{black}\\[6pt]
          \cafeteria\fontsize{20pt}{20pt}\selectfont\color{epnum}(S-17)\color{black}
          \\[18pt]
          % \flaregothic\fontsize{14pt}{14pt}\selectfont
          % \color{epattr}BIENVENUE, \uline{\color{Banpire} Mel} !
          \flaregothic\fontsize{18pt}{18pt}\selectfont
          \color{holop}L'HOLOPROJET !\\
          \color{black}
        }
      }
       &
      % Dòng thứ nhất: tên tập tiếng Nhật
      {\fotskip\fontsize{18pt}{18pt}\selectfont \ruby{助}{たす}けて！クビにされる！
      }                       \\[6pt]
      % Dòng thứ hai: tên tập tiếng Anh
       & {\cafeteria\fontsize{18pt}{18pt}\selectfont "Bloom," - Declaration Misleading}     \\[4pt]
      % Dòng thứ ba: tên tập tiếng Việt
       & {\vnmsans\fontsize{18pt}{18pt}\selectfont Bản kế hoạch tai hại}                     \\[8pt]
      % Dòng thứ tư: Nhân vật xuất hiện
       & {\caviar{\fontsize{14pt}{14pt}\selectfont Nhân vật: [\emp{kuon1} \emp{ryuuhou1} \emp{achan}] (Truyện phụ) \emp{matsuri} \emp{aqua} \emp{mio} \emp{flare1} \emp{marine} \emp{kanata}}
        }
      \\[8pt]
      % Dòng cuối cùng: Thời gian diễn ra của tập (thường sẽ là ngày phát hành)
       & {\continuum\fontsize{16pt}{16pt}\selectfont Ngày phát hành: 14/02/2021}            \\[36pt]
    \end{tabular}
  \end{adjustwidth}
\end{table}
%=========================================TRUYỆN CHÍNH========================================
\color{black}

\vspace{20pt}

\txtbx[2]
{{\color{vol} \uline{Tóm tắt:}} Khi Kuon nghỉ Tết, Mio, Flare, Matsuri, Kanata, Aqua và Marine tìm thấy một tờ giấy bí ẩn với những dòng chữ bị bôi đen. Họ hiểu nhầm đó là thông báo sa thải và hoảng loạn đến mức tổ chức một buổi ``huấn luyện idol\dots'' đầy những trò quái đản: trượt patin, bịt mắt sờ mặt, bánh mù tạt, và cả ném hào quang của Kanata. Cuối cùng, A-chan giải thích đó chỉ là bản kế hoạch ăn kiêng của một nhân viên, và họ được thông báo về dự án concert sắp tới. Cuộc gọi video với Kuon bị gián đoạn bởi sự xuất hiện bất ngờ của Aloe, khiến Kuon phải cúp máy vội vã.
}

Để chuẩn bị cho buổi live concert online ``Bloom,'' một sự kiện đặc biệt của \hll, nơi các thành viên có thể phô diễn tài năng của mình trên sân khấu ảo trước hàng triệu khán giả - ai nấy đều đang gấp rút luyện tập. Những ngày này, văn phòng trở nên nhộn nhịp hơn bao giờ hết, tiếng nhạc vang lên từ các phòng tập, tiếng bước chân dồn dập trên hành lang, và những cuộc họp khẩn cấp diễn ra liên tục. Nhưng giống như mọi lần, không khí chuẩn bị chẳng bao giờ có thể yên bình như mong đợi.

Sự bát nháo trong nội bộ trở thành điều không thể tránh khỏi. Người thì quên lịch tập, phải chạy vội đến phòng thu khi buổi tập đã bắt đầu được mười lăm phút. Kẻ thì chạy khắp nơi tìm trang phục biểu diễn, lục tung cả kho đồ mà chẳng thấy đâu. Có người còn lỡ tay làm đổ nước lên thiết bị thu âm, khiến bộ phận kỹ thuật suýt phát khóc vì phải chạy chữa trong đêm. Trong khi đó, một số thành viên lại mải mê chỉnh sửa lời bài hát hoặc tập vũ đạo đến mức quên cả ăn uống, để rồi khi kiệt sức thì mới nhận ra sai lầm, ngã vật ra sàn với cái bụng kêu réo.

Các cô gái thần tượng - như Mio, Flare, Matsuri\dots\ cũng không phải ngoại lệ. Họ đang vật lộn với những buổi tập kéo dài, những lần sửa lại vũ đạo, và những bản nhạc cứ thay đổi vào phút chót.

Và điều đáng nói nhất là họ đang không có một người quản lý tạm thời thay thế.

Hikawa Kuon - Kuon phụ trách điều hành chung - đã xin phép nghỉ về quê ăn Tết Nguyên Đán từ một tuần trước, ngay sau khi đã giải quyết hậu quả từ nàng Sư tử trắng, nàng Công chúa, nàng Ác ma và chính nàng Thuyền trưởng. Quyết định này không hề nằm ngoài dự đoán, bởi cậu đã làm điều tương tự vào năm ngoái, và cũng như lần trước, YAGOO đã dễ dàng chấp thuận.

Nhưng điều không ai ngờ tới là, chính vì quyết định này mà cậu đã vô tình chạm trán với một người mà đáng ra cậu không thể gặp - một người thuộc về thế hệ đã từng có năm thành viên nhưng giờ đây chỉ còn bốn. Cuộc gặp gỡ đó sẽ còn thay đổi rất nhiều thứ, nhưng đó lại là một câu chuyện khác, sẽ được kể ở Quyển D: Ký túc xá nhà Hikawa.

Còn hiện tại, câu chuyện vẫn tiếp tục ở đây - nơi các thành viên \hll\ đang đối mặt với một vấn đề hoàn toàn khác. Một vấn đề không liên quan đến những buổi luyện tập căng thẳng hay sự vắng mặt của quản lý, mà là một tờ giấy kỳ lạ. Một tờ giấy mà chẳng ai biết nó đến từ đâu, nhưng chắc chắn rằng nội dung trên đó có khả năng làm chấn động cả công ty. Nó nằm đó, vô tội trên bàn họp, như một quả bom hẹn giờ đang chờ được kích hoạt.

\il

Một tiếng nổ vang rền khắp văn phòng, rung chuyển cả những bức tường kiên cố. Ngay sau đó là hàng loạt tiếng hét thất thanh vang lên từ nhiều phía, hòa lẫn cùng tiếng kính vỡ và đồ đạc đổ ngã. Người thì chạy tán loạn, người thì ngã sõng soài xuống sàn nhà, tay ôm đầu với vẻ kinh hoàng. Đồ đạc bị hất tung tóe khắp nơi, giấy tờ bay lên trời rồi rơi xuống như những cánh hoa trắng trong một ngày giông bão.

Một người nào đó thốt lên trong bất lực, giọng như muốn vỡ tung qua đám khói bụi: ``\textbf{Mọi người à!! Bình tĩnh lại được không!?}''

Nhưng làm gì có ai còn giữ được bình tĩnh trong cái mớ hỗn độn này nữa. Tiếng khóc, tiếng la, tiếng chân chạy loạn xạ, tất cả hòa vào nhau tạo thành một bản giao hưởng của sự hỗn loạn.

Thật ra, đây cũng chẳng phải là lần đầu tiên văn phòng \hll\ rơi vào cảnh náo loạn như vậy, nhất là trong thời gian chuẩn bị cho một sự kiện quan trọng như ``Bloom,''. Với việc mọi người tất bật luyện tập, xử lý trang phục, chuẩn bị sân khấu, còn ban quản lý thì chạy đôn chạy đáo lo hậu cần, hỗn loạn gần như đã trở thành một phần không thể thiếu của công ty này. Người ta có thể nói rằng hỗn loạn là DNA của \hll.

Nhưng lần này thì có lẽ mọi chuyện đã đi hơi quá. Không ai rõ chuyện gì vừa mới xảy ra. Chỉ biết rằng, một vụ nổ đã xảy ra, mọi thứ thì rối tung lên, và giờ họ phải bắt tay vào dọn dẹp trước khi tình hình trở nên tệ hơn. Trong đám khói vẫn còn lảng vảng, những tiếng ho vang lên từ các góc phòng, cùng với những câu hỏi đầy hoang mang:

``Chuyện gì vừa xảy ra vậy?''

``Có ai bị thương không?''

``Ai đã gây ra vụ này!?''

Những câu hỏi thốt lên giữa đống hỗn loạn, nhưng chẳng có ai có thể đưa ra câu trả lời chính xác. Mọi người chỉ có thể đứng nhìn cảnh tượng đổ nát trước mắt, từng bước tiến về phía đống đổ nát để bắt đầu công cuộc dọn dẹp kéo dài. 

\ct

Khi những mảnh vụn đầu tiên được nhặt lên, ánh nắng xuyên qua cửa sổ lớn của văn phòng, chiếu rọi lên những mớ giấy tờ bay nhảy giữa không trung, dần dần xua tan đi đám khói mù mịt vừa mới bao trùm cả căn phòng. Từng chiếc bóng người hiện ra rõ hơn, cầm lấy chổi và xẻng, bắt đầu dọn dẹp từng góc nhỏ, từng mảnh vỡ bị bỏ lại sau vụ nổ.

``Lại thêm một ngày hỗn loạn nữa thôi\dots'' Flare thở dài dửng dưng, vừa cầm cây chổi vật vật quét mớ đổ nát trên sàn, mắt trông vô cảm trước cảnh tượng đồ đạc tung tóe khắp phòng.

Mio đứng bên cạnh, khuôn mặt đã đơ ra từ lâu khi cầm cây chổi của mình. Một tiếng ho khẽ bật lên vì hút quá nhiều bụi, rồi nàng cũng bực bội lẩm bẩm: ``Tại sao cái văn phòng này lại luôn rơi vào cảnh thế này nhỉ?''

Matsuri thì mệt lử ôm một hộp các-tông to đùng, chân bước lảo đảo tiến về phía hai người kia, miệng lầm bầm bất lực: ``Chúng ta được đào tạo để trở thành thần tượng chứ, mà giờ đây\dots\ sao lại trông như mấy người pha trò cho thiên hạ thế này\dots''

\img{7-11a}

Nàng Yêu tinh vàng chỉ cười khẩy, tay vẫn cầm chổi quét dồn đống mảnh vụn trên sàn: ``Công ty \hll\ nào mà chả bao giờ yên tĩnh, cũng đã quen với cảnh này từ lâu rồi.''

Nàng Sói đen cũng gật gù đồng tình, lắc đầu ngao ngán: ``Ngay cả câu nói đó cũng y hệt như Kuon-senpai từng nói hồi trước, không khác một chữ nào.''

Kanata đứng bên cạnh, vừa ôm cái hộp bìa cứng trong tay vừa bật cười trước lời than thở của hai đàn chị. Cô còn chưa kịp nói thêm gì thì Mio đã giục giã: ``Nhanh tay lên nào, chậm chạp thế này thì đến tối cũng chẳng dọn xong mớ hỗn độn này đâu.''

Hai người đàn em lười nhác chỉ đồng thanh đáp lời một cách miễn cưỡng: ``Vâng ạ\~{}''

Matsuri đang định đặt cái hộp xuống để tiếp tục công việc thì chợt thấy hai tờ giấy tuột ra từ khe hộp, rơi lộp bộp xuống sàn. Một tờ đã bị gấp nhiều lần, còn tờ kia vẫn phẳng phiu, chẳng khác nào vừa được đặt xuống chưa được bao lâu. Cô cúm người nhặt lên, ngẩn người nhìn những dòng chữ trên giấy một lúc, rồi gọi ba người kia lại với giọng đầy tò mò: ``Ê này, mấy người có biết đây là cái giấy gì không?''

Cô giơ tờ giấy lên, mặt đầy vẻ khó hiểu, đôi mày nhíu lại chặt: ``Chị vừa nhặt được nó lúc dọn dẹp này.''

Ba người còn lại lập tức kéo lại gần, những cái đầu tò mò đổ xô lại như một đàn chim sẻ đang xúm vào mẩu thức ăn. Mio cầm lấy một trong hai tờ giấy, nheo mắt nhìn kỹ từng dòng, đôi mắt cô tập trung cao độ như đang giải một bài toán hóc búa. Rồi cô chợt thốt lên, vẻ mặt như vừa nhận ra một thứ gì đó rất quen thuộc: ``Để xem nào\dots\ À, em biết rồi.''

``Có gì thế?'' Matsuri vội hỏi, giọng đầy háo hức muốn biết.

Mio không đáp ngay, chỉ im lặng thêm vài giây, mắt vẫn không rời tờ giấy, rồi cô nghiêng đầu trả lời với giọng như đang khẳng định một sự thật hiển nhiên: ``Cái này là giấy tờ của Kuon-senpai đấy.''

``Hả? Của anh ấy a?'' Kanata kêu lên kinh ngạc, đôi mắt mở to hết cỡ không dám tin, như thể vừa nghe thấy một tin tức chấn động cả công ty.

Nàng Sói đen gật đầu, giơ cao tờ giấy để ai cũng thấy rõ, ngón tay chỉ vào dòng chữ viết tay rất quen thuộc, nét chữ gọn gàng và rõ ràng: ``Đây là đơn xin nghỉ phép của anh ấy, chị nhận ra chữ viết của anh ấy liền.''

Không khí xung quanh bỗng im lặng một lúc, chỉ còn tiếng gió thổi lướt qua khe cửa và tiếng động nhỏ của ai đó đang dọn dẹp ở phòng bên cạnh.

``Nhưng sao đơn nghỉ phép của anh ấy lại lăn lộn ở đây được?'' nàng Thiên sứ nhíu mày, mắt không rời tờ giấy như thể nó vừa rơi từ trên trời xuống. Cô còn cầm lấy tờ giấy, lật đi lật lại như thể đang tìm thêm một dấu hiệu nào đó để giải thích tình hình.

Câu hỏi của Kanata khiến cả bốn cô gái đều đứng im, chẳng ai hiểu nổi lý do tờ giấy lại ở đây. Dù vậy, chỉ có một mình Mio trong số bốn người là biết được câu chuyện đằng sau - bởi cô chính là người cuối cùng gặp Kuon ở công ty từ một tuần trước.

``Chuyện là như thế này\dots'' nàng Sói đen thở dài, khoanh tay lại rồi bắt đầu kể vắn tắt câu chuyện, giọng cô như một người đang kể lại một ký ức vừa mới xảy ra.

\begin{center}
\uline{Hồi tưởng\\Một tuần trước.}
\end{center}

Sau khi giải quyết xong vụ rắc rối giữa nàng Botan, Towa và Luna, Kuon thở hắt ra một hơi dài, lấy lại bình tĩnh rồi đi thẳng đến phòng giám đốc, nơi Tanigo đang làm việc. Mục đích của cậu rất rõ ràng: xin phép nghỉ Tết Nguyên Đán để về quê, một truyền thống mà cậu luôn giữ gìn từ khi còn nhỏ.

Tất nhiên, cậu đã dự tính trước mọi khả năng. Nhưng giống như năm ngoái, ông chủ nhanh chóng đồng ý, không hề do dự. Ông hiểu rõ truyền thống của gia đình cậu, và cũng biết rằng cậu chưa bao giờ làm trái những quy tắc ấy, dù công việc có bận rộn đến đâu.

Sau khi nhận được sự chấp thuận, cậu quay lại văn phòng, trên tay cầm hai tờ giấy - một là bản gốc mang về, một là bản sao cậu để lại. Cậu tiến đến bàn làm việc của nàng Sói đen, đặt chúng xuống với một nụ cười nhẹ: ``Anh để lại bản sao tờ xin nghỉ phép của anh nha.''

Mio đang sắp xếp giấy tờ ở bàn làm việc, liếc nhìn cậu, chau mày với vẻ không hiểu: ``Anh có nhất thiết phải làm thế không?''

Cậu Tổng nhún vai, giọng điệu thoải mái như đang nói về một chuyện hiển nhiên: ``Cần chứ. Để cho mấy đứa khác vào mà nếu thắc mắc anh ở đâu thì cứ giơ tờ nghỉ phép đó cho mọi người xem.''

Nàng Sói đen khẽ thở dài, đưa tay xoa thái dương: ``Em thấy nó thừa thãi thế nào ấy\dots\ Em có thể nói một tiếng với họ là được mà.''

``Thừa còn hơn thiếu,'' cậu cười nhẹ, đặt tờ giấy xuống một góc bàn rồi vỗ nhẹ lên đó, như thể đó là một lời chào tạm biệt. ``Anh về Kawachi một chuyến, mấy đứa ở đây nhớ ngoan nha.''

``Thôi thì, em sẽ giữ nó vậy.'' Nàng Sói đen chống cằm, suy nghĩ một chút rồi hỏi với giọng đầy tò mò: ``Vậy khi nào anh về?''

Cậu Tổng ngẫm nghĩ một hồi như đang tính toán, ngón tay gõ nhẹ trên thành bàn, rồi đáp: ``Tầm mùng 8, mùng 9 Tết. Hôm nay là 26 Âm lịch, vậy là khoảng\dots''

Cô nàng nhẩm tính trong đầu rồi nhanh chóng kết luận, giọng cô như một phép toán đã được giải: ``12 đến 13 ngày. Tức là sau buổi concert của bọn em à?''

Cậu Tổng gãi đầu, vẻ hơi áy náy, đôi mắt cậu có chút hối lỗi: ``Ừ\dots\ Xin lỗi mấy đứa nha. Anh cũng muốn giúp mấy đứa vào buổi live sắp tới, nhưng ở nhà anh, gia đình luôn ưu tiên hơn so với công việc. Giám đốc cho anh nghỉ về quê ăn Tết đã là đãi ngộ rồi.''

Cậu cười nhẹ, nhìn nàng Sói đen đầy tin tưởng, giọng cậu như một lời gửi gắm: ``Anh nghĩ là mấy đứa sẽ đàng hoàng lúc anh về thôi mà, không phải lo cho anh đâu.''

Hôm đó, mặc dù trời mưa nặng hạt, Kuon vẫn nhanh chóng thu xếp hành lý, vội vã bắt xe buýt về thành phố Kawachi, để lại phía sau một văn phòng đang chuẩn bị cho sự kiện lớn nhất năm.

\begin{center}
\uline{Kết thúc hồi tưởng.}
\end{center}

Flare gật đầu từ từ, giọng cô dường như đã thấu hiểu vấn đề: ``Thì ra là vậy.''

Matsuri rũ người xuống ghế, chống cằm lên tay, tặc lưỡi với vẻ mặt quen thuộc khi đối mặt với những nét văn hóa lạ lẫm: ``Đúng là văn hóa ăn Tết Âm lịch của gia đình anh ấy khác biệt thật.''

Kanata thở dài, lắc đầu trong nỗi tiếc nuối: ``Thiệt là đáng tiếc nhỉ.''

Bốn cô gái ngồi im lặng một lát, ánh mắt trao đổi nhau rồi lại đắm chìm vào tờ đơn nghỉ phép của cậu Tổng. Cả nhóm đều cảm thấy hụt hẫng khi người quan trọng nhất - chỉ đứng sau YAGOO và A-chan - lại vắng mặt đúng vào thời điểm then chốt này, như một mảnh ghép quan trọng bị thiếu trong bức tranh đang hoàn thiện.

``\dots À, còn một tờ giấy nữa này,'' nàng Thiên sứ chỉ vào tờ giấy còn lại nằm trong tay Mio, đôi mắt long lanh nét tò mò.

Nàng Cổ động lật tờ giấy nằm phía sau đơn nghỉ phép lên. Trên mặt giấy in lớn dòng tiêu đề: ``Báo cáo và danh sách thay thế liên quan đến việc mở rộng và tác động trong năm nay\dots'', trông chẳng khác nào một tài liệu cấp cao từ ban lãnh đạo công ty.

Nàng Sói đen nhíu mày bối rối, đưa tay lấy tờ giấy xem kỹ hơn: ``Ủa, đây là cái gì thế?''

Cả bốn người cùng đăm đắm vào nội dung, nhưng phần lớn chữ đã bị bôi đen hoặc che mờ bằng những vệt mực dày đặc, y hệt một tài liệu mật đã trải qua quá trình kiểm duyệt gắt gao: ``Việc mở rộng [REDACTED] quá nhanh, đạt mức 100 [REDACTED] đã gây ra những tác động tiêu cực\dots''

Flare suy đoán, giọng tràn đầy lo âu: ``Cái này là đang nói về chúng ta à?''

Vừa dứt câu, nàng Thiên sứ mở to mắt, như vừa đọc phải một tin tức động trời. Nàng Yêu tinh vàng chỉ vào một đoạn trong báo cáo, ngón tay run rẩy: ``Đây này, con số 100 kia có lẽ là nói về việc chúng ta vượt mốc 100 triệu lượt xem, kiểu như vậy đấy.''

\img{7-11c}

Mio tiếp tục đọc, khuôn mặt dần chuyển sang màu nghiêm trọng, đôi lông mày xếch lại chặt chẽ: ``[REDACTED] và [REDACTED] thì chỉ khiến làm trò cười, điều đó cho thấy chúng ta đã đi lệch rất xa so với mục tiêu ban đầu.''

``Nghe giống như người viết đang phàn nàn rằng bọn mình cứ mãi làm trò hề quá phải không?'' Matsuri nhận xét, giọng đầy nỗi bất an.

Kanata tự nhiên cảm thấy lo lắng, bàn tay cô siết chặt lại: ``\dots Vậy đoạn cần loại bỏ ở đây là cái gì thế nhỉ?''

Cô đột ngột giật mình khi bắt gặp một câu trong báo cáo, mắt mở to trong kinh hoàng: ``Do đã bị xếp hạng D theo [REDACTED], chúng ta nên loại bỏ [REDACTED].''

Nàng Thiên sứ hét lên thất thanh, giọng như một tiếng sét đánh giữa trời trong xanh: ``Không lẽ là\dots\ \textbf{bị sa thải!?}''

Mio giật bắn mình, toàn thân cứng đờ như một bức tượng đá khi nghe hai từ mà không một ai muốn nghe nhất. Nàng Cổ động vội vàng trấn an, giọng đầy khẩn trương: ``Không thể nào đâu! Em đã là một idol chính thống mà!''

Nhưng Flare lại thốt ra một câu khiến mọi hy vọng bị dập tắt hoàn toàn, giọng như nhát chém cuối cùng: ``Nhưng trong giấy này có ghi 'không thể để tình trạng này tiếp diễn' còn gì.''

Nàng Cổ động lập tức câm lặng, miệng há hốc nhưng không thể phát ra được âm thanh nào. Trong nháy mắt, thế giới xung quanh họ như sụp đổ, những tia hy vọng mong manh bị dập tắt hoàn toàn. Cả ba người cùng lúc thốt lên những điều tiếc nuối chưa kịp hoàn thành, giọng như những tiếng kêu than cất lên từ sâu thẳm trái tim.

Và thế là\dots\ Đối với bốn cô gái - những thần tượng, những đứa con của gia đình Tam Giác Xanh - mọi thứ dường như đã đi đến hồi kết. Hoảng loạn trở thành phản ứng duy nhất mà họ có thể đưa ra.

Mio ôm đầu ngây người, hét lên trong kinh hoàng: ``\textbf{Waaaah!!}''

Matsuri khuỵu xuống sàn, giọng đầy bất lực và đau xót: ``\textbf{Chị còn chưa có cơ hội được nhìn đồ lót của Lamy-chan mà!}''

Còn nàng Thiên sứ thì ôm trán trong tuyệt vọng: ``\textbf{Em còn chưa trả hết khoản nợ tiền máy tính của mình nữa!}''

Riêng mình Flare vẫn đứng nguyên tại chỗ, vẻ mặt tròn xoe khó hiểu trước cảnh ba người kia phát cuồng, đầu cô nghiêng sang một bên như thể đang cố gắng giải mã một ngôn ngữ xa lạ mà mình chưa từng nghe thấy.

\img{7-11d}

Cô vội vàng lên tiếng trấn an, giọng đầy bối rối: ``Chắc chắn phải có sự nhầm lẫn ở đâu đó rồi\dots\ Sao chúng ta không thử liên lạc hỏi Kuon-senpai xem sao? Anh ấy là Tổng quản lý của mình mà, làm sao có chuyện không biết!''

Nhưng nàng Sói đen chỉ lắc đầu, giọng nói cạn kiệt hy vọng như một người đã đầu hàng trước số phận: ``Vô ích thôi. Dù có gọi bao nhiêu lần đi nữa thì anh ấy cũng sẽ không bao giờ nhấc máy đâu\dots''

Nàng Yêu tinh vàng nhíu mày thắc mắc: ``Sao lại có thể thế được?''

Tiền bối của cô thở dài nặng nề, giọng đắng chát: ``Bây giờ anh ấy đang ở nhà nghỉ Tết. Anh ấy đã để lại tất cả các thiết bị liên lạc với chúng ta ở công ty, chỉ mang theo máy có thể gọi được cho A-chan hoặc YAGOO thôi.''

``Chị cũng đã thử gọi cho anh ấy mấy ngày trước rồi mà cuối cùng cũng chẳng có ai nhấc máy\dots'' nàng Cổ động ngồi cạnh thêm vào câu chuyện, đầu vẫn cúi gằm xuống, giọng đầy thất vọng.

Nàng Yêu tinh vàng sững người, mắt mở to hết cỡ: ``Vậy nghĩa là\dots''

Nàng Sói đen gật đầu, đưa ra câu kết luận như một bản án cuối cùng: ``Dù có gọi hay nhắn tin thì cũng chẳng nhận được hồi âm từ anh ấy đâu.''

Không chịu thua, Flare tự mình bấm số gọi cho Kuon, nhưng chỉ nhận được câu trả lời lạnh lùng từ tổng đài điện thoại: ``Thuê bao quý khách vừa liên hệ hiện không thể kết nối. Xin quý khách vui lòng thử lại sau.'' Cô thử gửi một tin nhắn, màn hình chỉ hiển thị trạng thái ``đã gửi\dots'' chứ không chuyển thành ``đã xem\dots'', như một lá thư bị bỏ rơi giữa không trung vô tận.

Lúc này, cô mới thực sự hiểu rõ tình hình. Cô đứng lặng người trong cơn sốc, nhìn ba người kia kêu than trong tuyệt vọng với những lý do quá đỗi\dots\ đặc biệt. Cô không kiềm được thở dài trong bất lực: ``Mọi người có thể bình tĩnh lại một chút được không?''

Ngay lập tức, nàng Cổ động bật dậy, hét lên đầy kinh hoàng: ``\textbf{Làm sao mà bình tĩnh được chứ!? Ngày hôm nay có thể là ngày cuối cùng chúng ta được ở cùng nhau mà!!}''

Nàng Sói đen vẫn ôm đầu, than thở rên rỉ: ``Trời ơi, chuyện gì đang xảy ra với bọn mình vậy\dots''

Còn nàng Thiên sứ thì nước mắt đã lưng tròng, giọng nhỏ yếu như một đứa trẻ bị bỏ rơi: ``Kuon-senpai thực sự đã bỏ rơi chúng ta rồi\dots''

Nàng Yêu tinh vàng hoàn toàn bối rối trước cảnh tượng trước mắt. Một bản báo cáo bí ẩn, một tờ đơn nghỉ phép lạc lõng không rõ nguyên nhân, và một nhóm idol đang hoảng loạn vì những lý do khó lòng lý giải. Lẽ nào mọi chuyện lại có thể nghiêm trọng đến thế sao?

Cô lật ngược tờ giấy, cố gắng tìm ra bất kỳ thông tin nào còn sót lại. Nhưng những đoạn bị che mờ khiến cô không thể nào nắm bắt được toàn bộ nội dung, giống như một bức tranh bị che đi nửa phần bí mật.

``Này, mấy người có chắc là chúng ta không hiểu lầm gì đó đâu chứ? Có khi đây chỉ là một tài liệu nội bộ chưa hoàn chỉnh thôi?'' cô lên tiếng, đưa ra giả thuyết với hy vọng có thể xoa dịu bầu không khí căng thẳng đang bao trùm cả căn phòng.

Matsuri ngước mắt nhìn cô, vẻ mặt không khỏi nghi ngại: ``Em thực sự nghĩ vậy à?''

``\dots Ừ thì\dots'' cô đổ mồ hôi hột, cố gắng giải thích một cách hợp lý nhất: ``ít nhất cũng phải tìm cách xác nhận lại trước khi chúng ta hoảng loạn vô cớ chứ.''

Nàng Sói đen thở dài: ``Nhưng bây giờ chúng ta đâu còn cách nào để hỏi ai nữa\dots''

Không một ai trong số họ có thể liên lạc được với Kuon. Nếu muốn biết sự thật, họ phải tìm một người khác đủ thẩm quyền để giải thích về bản báo cáo này. Và đúng lúc đó, một tia hy vọng mong manh chợt lóe lên trong bóng tối tuyệt vọng.

Bất ngờ, trong đầu Kanata bỗng hiện ra một suy nghĩ, đôi mắt cô chợt sáng bừng lên: ``A-chan! Sao mình chẳng nghĩ ra sớm hơn? Chúng ta có thể hỏi chị ấy mà!''

Matsuri nghe vậy, đôi mắt cũng liền rực sáng, một tia hy vọng mới vụt sáng trong lòng: ``Đúng thật vậy! Chị ấy là trợ lý của giám đốc mà! Nếu có ai hiểu rõ chuyện này, thì chỉ có thể là A-chan mà thôi!''

Mio liên tục gật đầu, giọng nói chắc nịch như một mệnh lệnh không thể chối từ: ``Thế thì còn chần chờ gì nữa? Chúng ta mau đi tìm chị ấy thôi!''

Cả bốn người tức cuống cuồng cuộn gọn tờ báo cáo bí ẩn và đơn xin nghỉ phép lại, chẳng khác nào đang ôm một quả bom có thể nổ tung bất cứ lúc nào. Bàn tay run rẩy vì lo lắng, họ nhanh chóng nhét chồng giấy vào túi áo, không dám lãng phí thêm một giây nào cả. Cả nhóm hấp tấp xông ra khỏi phòng, để lại phía sau tiếng ghế bị xô lệch vang vọng. Họ chạy dọc theo hành lang văn phòng dài hun hút, mái tóc tung bay sau lưng, hơi thở dần dần trở nên gấp gáp, rít lên trong cổ họng bởi sự căng thẳng đến nghẹt thở.

Tiếng năm đôi giày dập dồn lên nền gạch lạnh lẽo, vang vọng khắp không gian trống vắng, tựa như hồi chuông cảnh báo đang thét lên về một bi kịch sắp sửa xảy ra. Mỗi bước chân đều chất chứa nỗi hoảng loạn đang dâng trào, họ chỉ có một suy nghĩ duy nhất trong đầu: phải tìm được A-chan, phải gặp được nàng Đeo kính để giải tỏa hết mọi nghi vấn, trước khi mọi chuyện trở nên quá muộn màng.

Vừa đến cửa phòng làm việc cá nhân của A-chan, nàng Thiên sứ đã vội vàng đẩy cửa vào mà không hề nghĩ đến việc gõ cửa, giọng nói vội vã như một lời kêu cứu khẩn cấp: ``\textbf{A-chan! Chị có ở trong phòng không!?}''

Nhưng căn phòng thì trống không. Không có bóng dáng của người trợ lý quen thuộc, chỉ còn lại chiếc bàn làm việc gọn gàng với vài chồng tài liệu được xếp ngay ngắn. Đèn trong phòng cũng đã tắt từ lâu, chứng tỏ chủ nhân của nó đã rời đi từ khá sớm, có lẽ là từ nhiều giờ trước.

``Lạ nhỉ? A-chan đi đâu rồi?'' nàng Cổ động thảng thốt lên, vừa nói vừa ngó nghiêng khắp căn phòng với khuôn mặt đầy hoang mang.

Nàng Sói đen bỗng chốc đứng sững người, ký ức ùa về như một cơn lốc. Cô vỗ mạnh vào trán, vẻ mặt bừng tỉnh như vừa nhớ ra một điều vô cùng quan trọng: ``Trời ơi, quên mất rồi\dots\ A-chan được công ty cho nghỉ mấy ngày để dưỡng sức mà!''

``\textbf{Hả!?}'' cả nhóm đồng thanh lên tiếng, đồng thanh ngỡ ngàng trước thông tin này, giọng nói hòa vào nhau tạo thành một bản hòa tấu của sự bối rối.

Cô thở dài, giải thích với giọng đầy sự thông cảm: ``Mấy hôm trước, A-chan làm việc xuyên đêm liên tục quá, suýt nữa thì ngất xỉu ngay tại văn phòng. YAGOO mới bắt chị ấy nghỉ ngơi vài ngày để hồi phục sức khỏe!''

``Vậy tức là\dots\ bây giờ chúng ta không thể hỏi chị ấy được nữa,'' đôi vai của Kanata liền rũ xuống, lộ rõ sự tuyệt vọng, giọng cô như một lời thừa nhận bất lực trước tình thế.

Matsuri thả người xuống ghế, khuôn mặt rầu rĩ của một người đã từ bỏ hy vọng: ``Thế thì phải làm sao bây giờ? Chúng ta chẳng hiểu rõ nội dung tờ báo cáo, lại cũng không thể liên lạc được với Kuon-senpai\dots''

``Có lẽ chúng ta phải tự mình tìm ra ý nghĩa của nó thôi,'' Mio nói, không còn che giấu được sự thất vọng trong giọng nói, như một tiếng thở dài não nề.

``Nhưng mà có quá nhiều chỗ bị che bôi đen kịt như thế này\dots'' Flare cầm tờ tài liệu nội bộ trên tay, liếc qua những đoạn bị mực đen che khuất, ``chẳng khác nào một người đang cố gắng nhìn xuyên qua một lớp sương mù dày đặc cả!''

Cả bốn người nhìn nhau với nỗi lo lắng hiện rõ trên mặt. Họ chẳng có bất kỳ đầu mối nào, càng suy nghĩ lại càng cảm thấy bất an vô cùng. Một bản báo cáo bí ẩn, một tờ đơn nghỉ phép bị thất lạc, và một cảm giác mơ hồ rằng có một điều gì đó lớn lao đang chờ đợi họ ở phía trước.

Đúng lúc đó\dots\ *CẠCH*

Tiếng cánh cửa bật mở vang lừng, lập tức kéo được sự chú ý của cả nhóm. Tất cả cùng quay đầu về phía cửa, chỉ thấy hai bóng dáng rón rén lẻn vào từ bên ngoài, y hệt hai con thú tò mò đang ngấp nghé bước vào hang của một loài khác.

Đó là Marine và Aqua, hai người thò đầu vào trong phòng mà chẳng hiểu chuyện gì đang xảy ra với đám bạn. Nàng Thuyền trưởng mở miệng trước, giọng đầy ngỡ ngàng: ``Ủa? Mấy người làm cái gì mà mặt mày nhìn thất thểu thế kia?''

Nàng Hầu gái nối lời ngay sau, giọng vừa như đang nhận xét một cảnh tượng lạ lùng: ``Sao mà trông ai cũng như sắp gục ngã, mất hết cả ý chí vậy?''

Bốn cô gái trong phòng đồng loạt hướng về hai người mới đến với cặp mắt trống rỗng, vô hồn, y như những kẻ đã đánh mất mọi thứ. Đôi mắt họ trông thật sự như những lỗ đen vũ trụ, hút cạn mọi tia sáng xung quanh. Bầu không khí im lặng đến đáng ngại khiến nàng Thuyền trưởng bất giác rùng mình, lùi lại một bước với vẻ mặt hoảng sợ: ``C-Cái gì vậy!? Đừng có nhìn em bằng ánh mắt đó chứ!''

\img{7-11e}

Không một ai trong nhóm đáp lời. Họ chỉ im lặng trao cho hai người mới đến tờ báo cáo đầy những đoạn bị bôi đen, đồng thời thều thào bằng giọng tuyệt vọng như những linh hồn đang than van: ``Chúng ta\dots\ sắp bị mất việc rồi\dots''

``\textbf{Hả!!?}'' / ``\textbf{Sao!?}''

Cả hai người vội vàng nhận lấy tờ giấy, đọc lướt từ đầu đến cuối, đôi mắt mở toang ra không thể tin được. Khi vừa hiểu được bề mặt nội dung của tờ báo cáo, họ cùng nhau hét lên trong kinh ngạc: ``\textbf{Cái gì cơ!? Bị đuổi việc á!?}''

``Nhưng mà nhưng mà-!'' Marine nói lắp bắp, mắt tròn xoe không thể tin vào những dòng chữ mình vừa đọc, bàn tay run rẩy nắm chặt tờ giấy: ``Em vốn đã là một idol trong sáng và tốt bụng từ trước đến nay mà!''

Lập tức, Aqua chằm chằm nhìn cô với ánh mắt đầy nghi ngờ, đôi mắt híp lại như một thám tử đang điều tra một vụ án bí ẩn: ``Em nói thật hả?'' Những người còn lại cũng cùng ngước nhìn nàng Thuyền trưởng với vẻ khó tin, như thể vừa nghe thấy câu chuyện tầm phào nhất trong năm.

\img{7-11f}

Ngay sau câu nói đó, nàng Thuyền trưởng đứng đơ người tại chỗ, như một bức tượng bị đóng băng giữa dòng thời gian. Cô gắng sức tìm lời để bào chữa, nhưng cuối cùng cũng chẳng nói được lời nào, chỉ biết quay mặt đi ho khan một tiếng để lảng tránh chủ đề, khuôn mặt đỏ bừng lên vì xấu hổ.

Matsuri thở dài một hơi, vỗ tay một cái thật mạnh để vực dậy tinh thần của mọi người. Cô khoanh tay, nhìn cả nhóm với vẻ mặt quả quyết của một vị tướng chuẩn bị bước vào trận chiến sinh tử: ``Chị biết bây giờ mọi người vẫn đang hoảng loạn lắm, nhưng vẫn chưa quá muộn đâu! Nếu hiện tại chúng ta chưa đạt được tiêu chuẩn của một idol, thì chỉ còn một cách duy nhất\dots''

``\dots Cách gì ạ?'' tất cả mọi người cùng hỏi một lượt, giọng nói hòa vào nhau như một bản hợp xướng xen lẫn giữa tò mò và một tia hy vọng mong manh.

``Chúng ta phải cùng nhau luyện tập, để nâng cao khả năng làm idol của mình thôi!''

Cả nhóm đứng ngẩn người trước lời tuyên bố đầy quyết tâm ấy, trên khuôn mặt ai cũng hiện lên vẻ vừa ngạc nhiên vừa nghi ngờ. Liệu rằng kế hoạch này có thực sự giúp họ thoát khỏi nguy cơ bị loại bỏ khỏi công ty hay không, đó là câu hỏi mà chỉ có thời gian mới có thể mang lại câu trả lời chính xác.

\ct

Với tư cách là ``huấn luyện viên\dots'' tự phong, Matsuri vừa mang về năm chiếc túi màu vàng rực rỡ, bên trong đựng toàn bộ đồ dùng cho buổi tập luyện sắp tới. Cô thả chúng xuống sàn với một tiếng ồn ào nặng nề, chẳng khác nào người vừa mang cả một kho báu về nhà.

Nhìn đống túi xa lạ trước mặt mà chẳng thể nào đoán được nàng Cổ động đang định xoay sở gì, Mio liền lên tiếng với chất giọng ngập tràn nghi vấn: ``Ê-tô\dots\ Cái đống này là sao thế nhỉ?''

\img{7-11g}

Với vẻ mặt đầy tự hào, nàng Cổ động đáp lời, giọng nói hào hứng chẳng khác nào người bán hàng đang nhiệt tình giới thiệu sản phẩm mới ra lò: ``Mấy chị vừa ra ngoài và mua hết mấy món đồ kinh điển của một idol đây!''

Ngay phía sau cô, bộ ba Aqua, Flare và Kanata đã tụ tập thành một nhóm nhỏ, những cái đầu tò mò đều cúi xuống để ngắm nhìn đống túi đồ kỳ lạ.

Nàng Hầu gái nghiêng người ngó vào trong một chiếc túi, lẩm bẩm với vẻ mặt đầy băn khoăn: ``Ở trong này có gì đâu mà bí ẩn thế nhỉ?''

Nàng Yêu tinh vàng nhún vai một cái, giọng nói thản nhiên như đang giải thích một chuyện hết sức đơn giản: ``Thì toàn đồ để tập luyện thành idol đó chứ sao. Em và Kanata đã cân nhắc kỹ lưỡng rồi, chọn mấy món mà bọn mình dùng được nhất.''

Nàng Thiên sứ cũng gật gù liên tục, giọng nói mang chút bí ẩn: ``Chọn đồ khó lắm chứ chẳng đùa đâu. Thôi thì cứ để đó xem Matsuri-senpai định làm gì với đống này nào.''

Khi cả nhóm vẫn còn đang ngơ ngác tò mò, Marine đứng bên cạnh bèn khoanh tay lại với vẻ hơi ngạo mạn, mạnh dạn lên tiếng tuyên bố như một người đã có vô số kinh nghiệm: ``Mặc vào thử đi đã, còn lại hỏi sau!'' Cô còn giơ ngón tay cái lên trời, như thể vừa nói ra một sự thật hiển nhiên nhất trần gian, với một biểu cảm vừa hài hước lại vừa đầy tự tin. 

Thấy vậy, nàng Sói đen bèn nhíu mày nhìn cô với vẻ cảnh giác: ``Sao chị nghe em nói mà thấy lòng này bất an thế không biết?'' Ngay trong đầu Mio chợt thoáng qua hình ảnh nàng Thuyền trưởng với bộ dạng ``chuẩn hải tặc\dots'', trông chẳng khác gì kẻ đang sẵn sàng bày ra một trò đùa rắc rối, với nụ cười gian xảo nở trên môi.

Matsuri vỗ tay một cái thật to, rồi hất tóc đầy tự tin với giọng điệu chẳng khác nào đạo diễn đang chỉ đạo một buổi diễn lớn: ``Mio-chan cứ đừng lo! Mấy đứa em đã tự chọn hết những món đồ phù hợp với mình rồi!'' Cô chống tay lên hông, quay sang nàng Sói đen rồi chỉ tay với vẻ đắc thắng, tựa như vị giám khảo đang sẵn sàng giới thiệu các thí sinh của mình: ``Mio-chan, em sẽ vào vai ban giám khảo nhé!''

Đôi tai sói của Mio bèn cụp xuống, cô cũng khoanh tay lại với ánh mắt đầy hoài nghi, giọng nói như người đã đoán trước được hết kết quả: ``Tôi cá là mấy người lại định đùa giỡn nhây nhể với nhau thôi!''

Dù có bị phản đối, nàng Cổ động vẫn chỉ vẫy tay xua đi mọi lo lắng với giọng đầy tin tưởng: ``Cứ yên tâm mà tin tưởng vào bọn chị đi\~{}''

Marine hào hứng tiến lên một bước, đặt tay lên ngực mình để tuyên bố với khí chất của một vị thuyền trưởng sắp ra lệnh xuất hành: ``Vậy để Senchou đây bắt đầu trước nhé!''

Cuối cùng, nàng Sói đen cũng chẳng thể làm gì khác hơn là thở dài, lẩm bẩm trong ngán ngẩm với vẻ mặt của một kẻ đã từ bỏ mọi hy vọng: ``\textit{Không biết bọn này có làm được trò gì ra hồn không nữa\dots}''

\ct

Buổi tập luyện bắt đầu với hình ảnh năm cô gái cùng nhau trượt giày patin. Những đôi giày liên tục lướt vút trên sàn nhà, tạo nên một bản giao hưởng đầy lộn xộn nhưng cũng tràn đầy sự cố gắng.

Mio chỉ ngồi chống cằm nhìn theo, khuôn mặt đầy hoài nghi, lòng thầm nghĩ: ``\textit{Đừng nói với tôi đây là bài kiểm tra dành cho idol đấy nhé\dots}''

Dù vậy, tất cả họ đều thực sự nghiêm túc - tất nhiên là theo cách rất riêng của mỗi người. Aqua và Kanata lảo đảo, bám chặt vào nhau để giữ thăng bằng, hai bàn tay đan xen như sợi dây cứu sinh giữa một biển khơi trơn trượt. Flare cũng cố gắng đứng vững, nhưng cơ thể vẫn lắc lư không ngừng như một con thuyền giữa cơn sóng lớn. Chỉ có Marine và Matsuri là dường như đã quen với trò này từ lâu, họ lướt đi một cách tự tin, chẳng khác nào những người đã nhiều năm trượt trên sân băng.

Khi mọi người đều trượt về phía trước, họ lần lượt vòng ra phía sau, để cho người đề xuất ra trò chơi này - nàng Thuyền trưởng - kiễng chân, giơ tay lên cao và giả vờ ngân một nốt cao. Giọng cô vang vọng khắp căn phòng, với vẻ mặt hoàn toàn tựa như một diva chuyên nghiệp.

``Đúng rồi đó! Một idol chân chính phải có bước di chuyển duyên dáng như vậy!'' nàng ``huấn luyện viên\dots'' Matsuri nhìn thấy mọi người dù còn vụng về nhưng vẫn cố gắng đồng bộ với nhau, ánh mắt tràn đầy tự hào như một người thầy đang chứng kiến học trò của mình tiến bộ.

Trong khi đó, Mio - người đã chuẩn bị sẵn một cặp kính râm và quàng chiếc khăn hồng chéo ngực, đóng vai một giám khảo khó tính - chỉ có thể thở dài bất lực: ``Tôi không nhớ có mục nào như vậy trong giáo trình đào tạo idol đâu\dots''

Sau vài phút nỗ lực, cả nhóm cùng quyết định thử diễn một màn trình diễn nhỏ. Họ xếp thành một hàng thẳng, dang rộng hai tay, bước chân đều theo nhịp nhạc, trông tựa như một nhóm vũ công nghiệp dư đang tập luyện cho một buổi diễn vô danh. Màn trình diễn kết thúc, năm người đứng thẳng tắp, khuôn mặt ai cũng tràn đầy hy vọng chờ đợi lời nhận xét.

\img{7-11h}

Và câu đầu tiên mà ``giám khảo\dots'' Sói đen thốt ra là: ``Tại sao ý tưởng về idol của mấy người lại lạc hậu đến thế!?'' Khuôn mặt cô không hề thay đổi sắc thái, giọng nói như một nhà phê bình kén chọn. ``Trông chẳng khác nào phong cách của thời Chiêu Hòa cả!''

Nghe vậy, Marine vừa tiếp tục trượt patin, vừa giải thích với nhiệt huyết của một nhà cách mạng âm nhạc: ``Văn hóa nhạc pop luôn luôn được làm mới! Luôn luôn như vậy mà--\dots''

Cô chưa kịp nói hết câu thì bỗng mất thăng bằng, trượt thẳng về phía Flare như một mũi tên bắn trượt khỏi mục tiêu: ``Khoan đã--!''

*RẦM!*

Cú va chạm khiến nàng Yêu tinh vàng chao đảo, còn nàng Thuyền trưởng thì kêu lên một tiếng ``Á\~{}'' đầy đáng ngờ, giọng nói như tiếng rên rỉ của kẻ đang tận hưởng một điều gì đó không đứng đắn. Cô nàng hải tặc dần lộ rõ bản tính ``hỏny\dots'' của mình, nhất là khi đôi mắt bắt đầu lấp lánh hình trái tim, tựa như một con thú săn mồi vừa tìm thấy con mồi ưng ý.

Nàng Thuyền trưởng lao tới với vẻ mặt sung sướng, giọng nói ngọt ngào như mật: ``Ufufu\~{} Matsuri-senpai à\~{}''

Nàng Cổ động nhìn thấy khuôn mặt của đàn em, hoảng sợ vội vàng quay đầu, vô tình túm lấy sợi dây thừng ở gần đó, như một người đang cố gắng trốn thoát khỏi một hiểm họa đang rượt đuổi.

*ỤP!*

Nàng Thuyền trưởng ôm chầm lấy đàn chị, đẩy ngã cô ấy xuống, kéo cả sợi dây kia theo mình, tạo thành một mớ hỗn độn chằng chịt trong phòng.

*TÁCH!*

Cả căn phòng bỗng chốc im lặng trong một giây ngắn ngủi, như thể thời gian đột nhiên ngừng trôi.

Mio, Aqua và Kanata đồng loạt ngước nhìn lên trần nhà. Một linh cảm chẳng lành ập đến, như làn gió lạnh lẽo chạy dọc sống lưng. Và rồi\dots

*BỐP!* *COONG!*

Một chiếc chậu nhôm từ đâu đó rơi thẳng xuống đầu nàng Thiên sứ, tạo ra một tiếng động vang dội. Nàng Sói đen giật mình vì tiếng động đó, cả đôi tai sói đều cụp xuống, trong khi nàng Hầu gái thì chẳng biết nên hoảng hốt hay bật cười trước cảnh tượng vừa xảy ra.

\begin{figure}[H]
\centering
\begin{subfigure}[t]{0.45\textwidth}
\centering
\includegraphics[width=8cm]{../../../../Image/Book7/7-11i1.png}
\caption{Marine vô tình va chạm với Flare, rồi lao thẳng về phía Matsuri\dots}
\end{subfigure}
\quad
\begin{subfigure}[t]{0.45\textwidth}
\centering
\includegraphics[width=8cm]{../../../../Image/Book7/7-11i2.png}
\caption{\dots khiến nàng Cổ động hoảng hốt túm lấy sợi dây thừng lơ lửng\dots}
\end{subfigure}
\quad
\begin{subfigure}[t]{0.45\textwidth}
\centering
\includegraphics[width=8cm]{../../../../Image/Book7/7-11i3.png}
\caption{\dots làm chiếc chậu nhôm rơi trúng đầu Kanata.}
\end{subfigure}
\caption{Tóm tắt chuỗi sự kiện hỗn loạn của nhóm năm cô gái.}
\end{figure}

Nàng Sói đen thở dài ngao ngán, giọng nói của một người đã chứng kiến quá nhiều chuyện bất ngờ: ``Được rồi\dots\ Tiếp theo!''

\ct

Sau đó là vô vàn những ý tưởng khác nối tiếp nhau - mà đúng hơn là một dãy những trò lạ lùng khó lòng nào rõ ràng liệu chúng có thực sự giúp họ hoàn thiện hơn với vai trò idol hay không, cứ như một chương trình hài kịch đang diễn ra mà chẳng thấy hồi kết bao giờ.

\il

Thử bắt đầu với trò chơi ``Thứ gì trong hộp?'', thế nhưng thay vì bỏ những vật thể bí ẩn vào bên trong, Marine và Matsuri lại quyết định bịt mắt mình lại, rồi dùng tay sờ khắp mặt của Kanata - cứ như hai người mù đang tập dượt để nhận dạng một bức tượng sừng sững.

\img{7-11j1}

Với vai trò là người ``nằm trong hộp\dots'', nàng Thiên sứ toát mồ hôi lạnh khi thấy hai bàn tay lảng vảng khắp khuôn mặt mình, giọng cô run lên như một tiếng kêu cứu tuyệt vọng: ``Khoan đã! Đó là lỗ mũi của em đấy mà!''

Flare chỉ biết khoanh tay đứng nhìn, đầu nghiêng sang một bên với vẻ mặt khó hiểu, chẳng sao nào thôi băn khoăn không biết trò này có liên quan gì đến công việc của một thần tượng. Còn Mio thì chỉ có thể xoa trán thở dài, giọng bất lực thốt ra một câu duy nhất: ``Tiếp theo!''

\il

Tiếp theo là trò ``Bốc bánh may rủi\dots'' theo đúng phong cách của Towa-sama - một thử thách đầy rủi ro, vì có một chiếc bánh được nhồi đầy mù tạt cay xé. Nàng Hầu gái, nàng Cổ động và nàng Sói đen đều may mắn an toàn, cả ba cùng thở phào nhẹ nhõm khi thoát khỏi cái bẫy. Riêng nàng Thiên sứ thì\dots

``!\$)*\$\#(\$\#!\$\~{}\~{}\~{}!!!'' cô nàng rên siết trong cơn đau, nước mắt lưng tròng giàn giụa, khuôn mặt đỏ bừng như lửa vì vị cay xé lưỡi. Cô ôm miệng chạy vòng khắp phòng, chẳng khác nào một con thú vừa bị thương nặng.

\img{7-11j2}

Dù thở phào mừng vì mình không lại dính bẫy mù tạt lần nữa\footnote{Xem lại {\flaregothic ÉPISODE 79}, phần {\abv Truyện phụ}, quyển số Sáu.}, nàng Sói đen chẳng hề tỏ ra chút thương xót, giọng như một vị giám khảo khó tính thốt lên: ``Tiếp theo!''

\il

Chẳng dừng lại ở đó, Marine và Aqua còn nghĩ ra một trò còn tàn nhẫn hơn: dùng chính vầng hào quang hình ngôi sao bốn cánh trên đầu nàng Thiên sứ để làm đồ chơi ném bắt, cứ như đang chơi tung hứng với một vật thể quý giá lắm.

\img{7-11j3}

Nàng Hầu gái cười tươi rói, quăng cái hào quang về phía nàng Thuyền trưởng như một cái đĩa bay, giọng đầy hào hứng: ``Nào Marine-senchou! Cứ bắt lấy đi!''

Nàng Thuyền trưởng giơ tay chờ sẵn, chộp lấy ``của quý'' của cô nàng thế hệ thứ tư này với vẻ mặt hớn hở, chẳng khác nào một kẻ săn mồi vừa bắt được con mồi ưng ý: ``Được rồi, để em giữ nó!''

Ở giữa, nàng Thiên sứ hốt hoảng chạy tới lui để giật lại thứ vốn thuộc về mình, giọng gần như khóc nức nở: ``\textbf{Không được! Trả lại cho em cái đó đi!!}''

\il

Và cứ thế, vô vàn trò ngốc nghếch cứ nối tiếp nhau không ngừng nghỉ: từ việc nhảy dây bằng chiếc băng đô của Flare, đến lúc Matsuri cố gắng khiến Aqya hát thật to trong khi đang ngậm nước, rồi Marine lại trình diễn màn ``biến thành idol\dots'' nhưng chỉ toàn những tư thế tạo dáng kỳ quặc hết mức. Tất cả cứ như một vở hài kịch đang diễn ra độc diễn trước một khán giả duy nhất trong căn phòng.

Ngay cả từ bên ngoài cửa sổ, ai cũng có thể nghe thấy đủ thứ âm thanh hỗn loạn vang lên: tiếng la hét thất thanh của nàng Thiên sứ, tiếng hô hào thúc giục của những kẻ đang quá đà hào hứng, tiếng cười khúc khích đầy tinh quái, và luôn xen lẫn là những tiếng thở dài ngao ngán của nàng Sói đen. Tất cả hòa quyện vào nhau thành một bản giao hưởng của sự điên rồ.

Cuối cùng, nàng Thiên sứ tuyệt vọng la lên, giọng như một đứa trẻ đang khóc nghẹn: ``\textbf{Dừng lại hết đi! Em nói nghiêm túc đấy, mấy chị đang vượt quá mọi giới hạn rồi!!}''

Nàng Thuyền trưởng cười nhếch mép, giọng đầy thách thức đáp lại: ``Câm mồm đi! Em nói xem, bây giờ em là gì?''

Nàng Thiên sứ mếu máo đáp, giọng nhỏ bé như một đứa trẻ vừa bị bắt nạt: ``Là\dots\ là idol ạ!''

``Và bọn ta đang ở đâu?''

``\textit{hololive\dots}''

Nàng Thuyền trưởng chỉ thẳng lên trời, giọng như một vị tiên tri đang tuyên bố chân lý lớn lao: ``Thế thì chúng ta có gì phải sợ nữa, đúng không nào!? \textbf{Đúng không nào!?}''

Nàng Yêu tinh vàng giật mình hoảng hốt, la lên để ngăn cản: ``Marine, dừng lại đi nào!''

Nàng Cổ động thì chống nạnh, giọng đầy quả quyết như một huấn luyện viên đang động viên học trò: ``Chúng ta có đang tập luyện để trở thành idol tốt hơn mà!''

Nàng Hầu gái cuống cuồng vung tay lên, giọng lo lắng thốt ra: ``Mọi người có thể lùi lại một chút được không!?''

Còn nàng Sói đen thì hoảng loạn hét lên như một tiếng kêu cứu: ``\textbf{Ai đó lại mau cầm lấy sợi dây thừng kia đi chứ!}''

Họ cứ thế cãi nhau, chất vấn nhau, rồi lại tiếp tục những trò kỳ quặc đó cho đến tận chiều hôm đó, cứ như một vở kịch vô tận đang diễn ra trên một sân khấu không có khán giả. Nhưng cuối cùng, chẳng có trò chơi nào trong số đó thực sự giúp họ cải thiện kỹ năng làm thần tượng như ban đầu mong muốn. Tất cả chỉ để lại tiếng cười rộn ràng của người xem nếu vô tình bắt gặp và những cơn đau đầu nhức nhối cho người duy nhất đứng nhìn tất cả.

\ct

Bóng hoàng hôn buông dần, những tia nắng cuối ngày len lỏi qua ô cửa sổ, tô điểm toàn bộ căn phòng bằng sắc cam ấm áp mà cũng phảng phất sự mệt mỏi sau một ngày dài.

Hết thảy mọi người đã kiệt quệ. Flare dựa lưng vào ghế sô-pha, đầu ngả ra sau, mắt thờ ơ dán vào trần nhà, từng hơi thở dài cứ thế trôi đi. Bởi một lý do khó hiểu, Marine bị trói tay, ngồi bệt dưới sàn với hơi thở hụt hẫng, chẳng khác nào một thủy thủ vừa vượt qua cơn bão biển khủng khiếp. Còn Kanata thì nằm bẹp dí không động đậy, linh hồn dường như đã rời xa thể xác, chỉ còn lại một thân thể vô hồn trên sàn gỗ.

Matsuri khuỵu xuống, đầu gục xuống phía trước, lẩm bẩm trong nỗi thất vọng bao trùm, giọng cô yếu ớt như một tiếng than thở từ sâu thẳm trái tim: ``Không ổn rồi\dots\ \textbf{Chúng ta đang tập luyện để trở thành diễn viên hài, chứ đâu phải idol!}''

\img{7-11k}

Nàng Thuyền trưởng rên rỉ, khuôn mặt tràn đầy sự bất lực và bi ai: ``Nếu cứ tiếp tục như thế này, chúng ta sẽ kiệt sức mà chết mất!''

Vừa dứt lời, cô cũng sụp xuống sàn, mắt nhắm chặt, đầu nghiêng sang một bên tựa một chiến binh đã gục ngã trên chiến trường. Nàng Thiên sứ chỉ có thể thở ra một tiếng thở dài chua chát, chẳng còn sức để phản ứng gì nữa, dường như cô đã từ bỏ mọi hy vọng.

Trên chiếc ghế sô-pha, Aqua và Mio cũng không khá hơn chút nào. Nàng Sói đen thậm chí còn úp mặt xuống gối, đôi tai cụp lại đầy chán nản, trông giống hệt một con sói vừa đánh mất con mồi sau bao ngày săn đuổi.

Không một ai mở lời. Sự im lặng bao bọc toàn bộ văn phòng, chỉ còn lại tiếng thở đều đặn của những con người đã cạn kiệt năng lượng. Dường như tất cả đều buông bỏ, chấp nhận số phận rằng họ có thể bị sa thải bất cứ lúc nào, như một kết quả không thể thay đổi.

Nhưng có một người ngoại lệ.

Chính là cô nàng Hầu gái mang biệt danh ``Baqua\dots'' - người đã suy ngẫm về mọi chuyện xảy ra trong ngày, cũng là người nhiệt thành nhất với buổi ``huấn luyện\dots'' này. Ngồi yên đó, đôi mắt tím long lanh của cô phản chiếu ánh đèn mờ ảo trong phòng, như hai viên pha lê lấp lánh giữa bóng tối. Cô nhìn vào lòng bàn tay mình, rồi bất giác siết chặt lại. Một cảm giác mơ hồ dần len lỏi trong tâm trí, một suy nghĩ mà cô luôn mang trong lòng nhưng chưa bao giờ nói ra thành lời.

Cuối cùng, để phá tan bầu không khí u sầu, cô cất giọng nhẹ nhàng nhưng đậm chất suy tư: ``Nhân tiện\dots\ cuối cùng thì điều gì tạo nên một idol nhỉ?''

Câu hỏi bất ngờ khiến cả phòng đột ngột im lặng, như một hòn đá rơi xuống mặt hồ phẳng lặng, tạo ra những vòng nước lan rộng khắp mặt nước. Mio khẽ nhấc đầu lên, ánh mắt trầm ngâm. Kanata liếc nhìn Marine, trong khi Flare nghiêng người về phía trước, hai bàn tay đan xen vào nhau.

Ba người đang nằm bẹp trên sàn lần lượt ngước nhìn cô, rồi đáp lại theo bản năng, nhưng trong giọng nói vẫn còn chút mơ hồ, dường như chính họ cũng chưa thực sự chắc chắn về câu trả lời của mình.

Matsuri là người đầu tiên lên tiếng, giọng có chút bỡ ngỡ: ``Ừm\dots\ Là khi chúng ta tỏa sáng rực rỡ trên sân khấu sao?''

Kanata tiếp lời, giọng cô nhẹ nhàng như một giấc mơ: ``Là khi ta đẹp và tài giỏi như một thiên thần?''

Marine nở một nụ cười yếu ớt, giọng đầy vẻ hài hước: ``Hay đơn giản chỉ là khi chúng ta luôn cố gắng hết sức thôi?''

Nàng Hầu gái im lặng, để những lời nói đó trôi qua như làn gió thoảng. Mọi ý kiến đều có lý, nhưng sâu trong tim, cô biết mình đang tìm kiếm một điều gì đó khác biệt, một điều gì đó không thể diễn tả bằng những từ ngữ đơn giản.

Cô cúi đầu, đặt hai bàn tay lên đầu gối, siết nhẹ rồi thả lỏng. Đôi vai nhỏ bé của cô run lên một chút trước khi hít một hơi thật sâu, dường như đang chuẩn bị cho một điều vô cùng quan trọng.

Rồi cuối cùng, cô lại hít một hơi sâu nữa và cất tiếng, giọng có chút nghẹn ngào nhưng vẫn tràn đầy sự kiên định: ``Đối với chị thì\dots''

\img{7-11l}

Cô ngừng lại một lát, như đang cân nhắc từng chữ, từng lời trước khi nói ra. Nàng Yêu tinh vàng hơi nhướng mày, còn nàng Thiên sứ thì quay hẳn người sang nhìn cô với ánh mắt tò mò.

Khi đã bình tĩnh trở lại, nàng Hầu gái tiếp tục, giọng nói dần vững vàng hơn: ``Chị nghĩ rằng, một idol là người có thể tỏa sáng và mang lại nụ cười cho khán giả, trong khi họ dồn hết tâm huyết cho những gì mình đam mê.''

Nàng Cổ động chớp mắt, dường như vừa chợt hiểu ra một điều gì đó. Cô đã từng nghe những lời tương tự, nhưng khi chúng được thốt ra từ miệng Aqua - người luôn vụng về nhưng chưa bao giờ từ bỏ - chúng bỗng trở nên sâu sắc và nặng nề hơn, như một lời thú nhận từ tận đáy lòng.

``Chị không quan tâm người khác nói rằng những gì mình làm chẳng giống một idol chút nào,'' cô tiếp tục, siết chặt vạt váy, các ngón tay cô nắm chặt vào nhau.

Marine hơi há miệng nhưng lại không nói nên lời. Cô biết nàng Hầu gái đã từng đấu tranh rất nhiều với danh xưng ``idol'', đặc biệt khi vô số người ngoài kia đặt ra những tiêu chuẩn khắt khe đến mức phi lý, như một bức tường vô hình khó lòng vượt qua.

``Đối với chị, cách mình tỏa sáng đơn giản chỉ là chơi game và trò chuyện cùng mọi người trên stream thôi,'' cô mỉm cười, nhưng nụ cười đó lại phảng phất chút suy tư, một chút tự hào xen lẫn sự khiêm tốn. Những giọt nước mắt rơi thẳng xuống đôi tay nhỏ nhắn của cô, lấp lánh dưới ánh đèn phòng.

Mio ngẩng đầu lên, ánh mắt thoáng rung động. Dù đôi khi có trêu chọc đàn chị vì sự vụng về, nhưng giờ đây cô hiểu rằng đó không phải là thiếu cố gắng, mà là cách riêng của nàng Hầu gái để kết nối với thế giới - một sự kết nối chân thành, không chút giả tạo.

Aqua nắm chặt hai bàn tay lại, tiếp tục dòng tâm sự của mình, giọng nói êm đềm như dòng chảy: ``Mình muốn làm những điều mình yêu thích, cùng với tất cả các fan và các thành viên của \hll.''

Cô dừng lại, đôi mắt ánh lên sự chân thành tuyệt đối, một ngọn lửa mà không ai có thể nghi ngờ. Cô khẽ ngước nhìn mọi người, mắt vẫn còn long lanh lệ, nhưng giờ đây toát lên vẻ kiên định vô cùng: ``Và đó mới chính là kiểu idol mà mình muốn trở thành.''

Sự im lặng lại bao trùm căn phòng, dày đặc như một tấm màn nhung.

Không ai nói một lời, nhưng ánh mắt của tất cả dần thay đổi. Không còn sự chán chường hay mệt mỏi, thay vào đó là sự thấu hiểu và đồng cảm, dường như họ vừa nhìn thấy một điều gì đó mới mẻ trong chính bản thân mình.

Lời bộc bạch của cô vang vọng trong căn phòng tĩnh lặng, như một ngọn đèn nhỏ thắp sáng bóng tối, một tia sáng le lói giữa đêm dài. Từng câu chữ đơn giản nhưng đầy chân thành khiến tất cả đều lắng nghe một cách nghiêm túc. Và dần dần, họ nhận ra: những gì cô vừa nói chính là sự thật - một sự thật mà họ luôn biết nhưng chưa bao giờ dám thừa nhận.

Flare khẽ tựa cằm lên tay, khóe môi nở một nụ cười nhẹ. Kanata chớp chớp mắt, như đang cố kìm nén những cảm xúc dâng trào trong lòng. Matsuri ngả người ra sau, mắt nhìn lên trần nhà, miệng khẽ cười vì vừa chợt hiểu ra một điều vô cùng quan trọng.

Marine nhắm mắt, rồi thở ra một hơi dài. Cô nở một nụ cười tinh nghịch nhưng có phần sâu lắng: ``Thật là\dots\ Akutan lại khiến cả phòng rơi vào không khí nghiêm túc thế này\dots''

Mio bật cười nhẹ, nhưng trong mắt vẫn còn đọng lại chút xúc động sâu sắc: ``Có lẽ Aqua nói đúng.''

Không ai phản bác. Bởi tất cả đều hiểu - idol không chỉ đơn thuần là đứng trên sân khấu, không phải chỉ có những hào quang chói lọi hay những bộ váy lấp lánh. Mỗi người đều có cách tỏa sáng riêng của mình, và con đường mà Aqua chọn hoàn toàn không sai.

Nàng Sói đen nhìn tiền bối một lúc lâu trước khi mỉm cười, dường như cô vừa tìm thấy câu trả lời cho chính bản thân mình. cô nắm lấy tay người đàn chị, nở một nụ cười dịu dàng: ``Mọi người ai cũng có định nghĩa về idol cho riêng mình. Chúng ta biết mình muốn trở thành một idol như thế nào để mang lại niềm vui lớn nhất cho khán giả.''

Một câu nói không cần to tiếng, nhưng đủ sức nặng để bất cứ ai nghe được cũng phải suy ngẫm rất lâu. Bởi đó là sự thật. Idol không phải là một khuôn mẫu cố định mà ai cũng phải theo - nó là hành trình mà mỗi người tự vẽ nên cho mình, một bức tranh không bao giờ lặp lại.

Matsuri nghiêng người về phía trước, ánh mắt bừng sáng như một tia lửa vừa bùng cháy: ``Đúng vậy! Việc phán xét và đuổi việc ai đó chỉ dựa trên những tiêu chuẩn hời hợt đó thật sự là bất công!''

Cô siết chặt tay, vẫn còn bức xúc về tình cảnh họ đang đối mặt, ngọn lửa phẫn nộ đang cháy trong lồng ngực. Nàng Thuyền trưởng cũng gật gù đồng tình, đặt một tay lên hông với giọng đầy quyết tâm: ``Vậy ngày mai chúng ta cùng đến nói chuyện với cấp trên nhé!''

Flare mỉm cười, ánh mắt tràn đầy tự tin mạnh mẽ: ``Tôi đồng ý! Nghe có vẻ đó là kế hoạch hoàn hảo cho chúng ta đúng không?''

Kanata siết chặt nắm tay, giọng cô như một lời thề: ``Nếu không thể nói chuyện với Kuon-senpai, chúng ta có thể gặp lại A-chan, thậm chí là YAGOO nữa!''

``Để mọi chuyện phải được giải quyết ổn thỏa!'' tất cả cùng giơ nắm đấm lên trời, ánh mắt tràn đầy quyết tâm cao độ, tựa một đội quân chuẩn bị bước vào trận chiến cuối cùng.

Bầu không khí chán chường và thất vọng lúc trước đã biến mất hoàn toàn. Ngọn lửa quyết tâm tìm ra sự thật bùng cháy trên gương mặt từng người, đôi mắt họ sáng bừng lên với hy vọng mới.

\img{7-11m}

Aqua nhìn mọi người với sự ngạc nhiên, đôi mắt tím tròn xoe như hai viên bi ve. Cô chưa bao giờ nghĩ rằng chỉ với những suy nghĩ chân thành của mình lại có thể tạo ra động lực lớn đến vậy. Những người chị em bên cạnh, họ thực sự tin tưởng vào điều này, và điều đó làm trái tim cô ấm áp vô cùng.

Cô khẽ cúi đầu, môi mím lại như đang suy nghĩ điều gì đó. Nhưng khi ngẩng lên lần nữa, trên gương mặt nhỏ nhắn đã nở một nụ cười nhẹ. Một nụ cười không cần lời nói, nhưng ai cũng hiểu rằng cô đã đồng ý, đã sẵn sàng cho cuộc đối đầu sắp tới.

Mio nhìn quanh, thấy ai cũng đã sẵn sàng. Cô mỉm cười, ngẩng đầu lên đầy tự tin, giọng cô như một tiếng hiệu lệnh: ``Vậy là tất cả đều đồng tình với ý kiến này chứ?''

Tất cả cùng gật đầu đồng tình, những cái gật đầu tràn đầy sự quyết tâm.

``Vậy thì chúng ta hãy cùng đòi lại công bằng cho chính mình! \textbf{Cố lên!}''

Flare, Kanata, Marine (giơ cả hai tay), Matsuri và Aqua cùng nhiệt liệt hưởng ứng với giọng nói vang dội: ``\textbf{Cố lên!}''

Họ không còn là những idol bị động, luôn bị động với mọi tình huống. Họ đã sẵn sàng đứng lên bảo vệ con đường mình đã chọn, một con đường mà không ai có thể định nghĩa thay họ.

\begin{center}
\uline{Sáng hôm sau.}
\end{center}

Không khí buổi sáng vẫn còn vương chút se lạnh, những tia nắng sớm mỏng manh xuyên qua các ô cửa kính của công ty, vẽ nên những vệt sáng dài trên sàn nhà gỗ. Từ ngoài đường phố, tiếng còi xe và bước chân vội vã của mọi người đang chuẩn bị cho ngày mới đều vang vọng lại, nhưng sáu cô gái chẳng có tâm trạng nào để thưởng thức bữa bình thường ấy. Họ bước vào công ty với gương mặt dồn đầy quyết tâm, tựa như một đội quân đang chuẩn bị bước vào trận chiến cuối cùng.

Mio đi dẫn đầu, lần lượt phía sau là Flare, Kanata, Matsuri, và lần này cả nàng Aqua và Marine đều đồng hành. Sáu người cùng nhau tiến về phía trước với từng bước chân hừng hỡi khí thế, như thể đang lao thẳng vào một cuộc đối đầu không còn cách nào tránh được.

Vẫn như mọi khi, nàng Cổ động là người chủ động nhất. Cô đẩy mạnh cánh cửa văn phòng với một lực mạnh đến mức cánh cửa bật ra và va chạm vào tường, giọng cô vang vọng khắp hành lang như tiếng sấm: ``\textbf{Quản lý có ở đây không ạ!?}''

\img{7-11n}

Đáp lại lời kêu gọi đó chỉ có một sự im lặng kỳ lạ, đến mức có cảm giác căn phòng đã bị bỏ trống từ rất lâu rồi.

Không một ai trong nhóm thấy được nhân viên nào trong phòng cả, chẳng có tiếng gõ phím nào, cũng không có tiếng trò chuyện nào vang lên. Không gian yên tĩnh đến nỗi có thể nghe rõ cả nhịp thở của chính mình, một sự bình yên khác thường hoàn toàn trái ngược với không khí náo nhiệt thường ngày của văn phòng.

Nàng Thiên sứ nheo mắt nhìn quanh, giọng lẩm bẩm đầy lo lắng: ``Em không thấy ai trong phòng cả.''

Nàng Yêu tinh vàng khoanh tay, giọng nghi hoặc như một nhà thám tử đang điều tra vụ việc: ``Có lẽ họ đang sao nhãng quá? Hay là\dots\ không một ai tới đây làm việc? Hôm nay là thứ Hai mà\dots''

Nàng Hầu gái, vốn ít khi chủ động, bỗng nhiên chỉ tay về chiếc bàn làm việc gần cửa sổ, giọng đầy ngạc nhiên: ``Không đâu. A-chan đang ở kia kìa.''

Tất cả mọi người lập tức dõi theo hướng tay chỉ của cô. Ở một góc văn phòng, một cô gái đeo kính đang tập trung hoàn toàn vào màn hình, ngón tay gõ phím liên tục với tốc độ của người đang chạy đua với thời gian. Cô ấy dường như đang chuẩn bị một bài thuyết trình trực tuyến, những dòng chữ và biểu đồ liên tục trượt qua màn hình như một dòng sông chảy xiết.

\img{7-11o}

Cả nhóm cùng nhau thở phào nhẹ nhõm, đôi vai ai cũng hạ xuống một cách rõ rệt. Ít nhất thì có vẻ A-chan vẫn tới đang làm việc bình thường, một tia sáng nhỏ nhoi xuyên qua mây mù lo lắng của họ.

Vừa nghe thấy tiếng động, nàng Đeo kính lập tức quay đầu lại, nở một nụ cười nhẹ nhàng với giọng điềm tĩnh: ``Ồ, chào buổi sáng mọi người. Sáng sớm ra có chuyện gì mà gấp gáp thế?''

Nhưng trước khi kịp nói thêm điều gì, Aqua đã chạy thẳng đến chỗ cô, gương mặt đầy nghiêm nghị như một chiến binh sắp bước vào trận chiến, giọng nói vang lên đầy khẩn trương: ``A-chan à!''

Nàng Đeo kính giật mình trước sự nhiệt tình bất ngờ này, cô khẽ lùi lại một chút với vẻ mặt không hiểu chuyện gì đang xảy ra: ``S-Sao vậy?''

Nàng Hầu gái hít một hơi thật sâu rồi nói, giọng nói rõ ràng và đầy quyết tâm như đang đọc một bản tuyên ngôn: ``Có thể bọn em vẫn còn rất xa mới có thể trở thành một idol hoàn hảo\dots''

Mio bước lên, từ từ tiếp lời với giọng đầy tự tin: ``\dots nhưng chỉ có bọn em mới là người quyết định được bọn em sẽ trở thành một idol như thế nào.''

Flare gật đầu, đôi mắt sắc bén toát lên sự tự tin, giọng nói như một lời khẳng định chắc chắn: ``Cùng với các fan hâm mộ, bọn em sẽ từng bước tiến gần hơn đến mục tiêu của mình.''

Matsuri siết chặt bàn tay, giọng đầy bức xúc như một người đang đòi lại công lý: ``Và bọn em nghĩ rằng thật bất công nếu có ai đó đánh giá và sa thải bọn em chỉ vì những lý do đó ạ.''

Marine bước lên, hai tay ôm chặt ngực, giọng đầy bi thương như đang đóng vai trong một vở kịch buồn, đôi mắt long lanh như đang xúc động: ``Bọn em muốn chị cân nhắc lại những đánh giá và kế hoạch sa thải này ạ. Nếu không thì\dots''

Cô dừng lại một giây, rồi đột ngột đổi giọng thành tiếng kêu đau khổ, quỳ xuống một cách thảm thương như một thuyền trưởng đang chứng kiến con tàu của mình chìm dần xuống biển, giọng nói như một tiếng thét tuôn từ đáy lòng: ``Thuyền của em\dots\ Không, giấc mơ của em sẽ chìm nghỉm mất!''

Kanata theo phản xạ cũng bật lên một câu đầy lo lắng, giọng như tiếng kêu cứu tuyệt vọng: ``Em còn nợ phải trả nữa! Nếu bọn em bị đuổi, em đào đâu ra tiền để trả đây\dots?''

A-chan đứng yên, dần dần ổn định cảm xúc sau dòng lời dồn dập mà sáu cô gái vừa xướng lên. Cô chỉnh nhẹ gọng kính, ánh mắt vẫn còn ngẩn ngơ chưa kịp nắm bắt hết tình hình, nhưng nét cười nhẹ đã khẽ nở trên đôi môi. Cuối cùng, cô thở ra thật nhẹ rồi nói với giọng bình thản: ``Trước hết, chị muốn nói rằng công ty hoàn toàn không có ý định sa thải ai cả, kể cả Hikawa-san cũng vậy.''

Cả nhóm đứng đơ như tượng, mọi chuyển động đều bị đóng băng trong nháy mắt.

Nàng Đeo kính tiếp tục, giọng nói vừa điềm đạm vừa ấm áp: ``Bọn chị cũng nghĩ việc các em phát triển theo phong cách riêng của mình, trở thành những idol độc đáo của bản thân, là một điều vô cùng quý giá.''

Nàng Hầu gái nhíu mày, chớp mắt mấy lần trước khi cất lời, như thể đang kiểm tra lại xem mình có nghe nhầm không, giọng cô đầy bối rối: ``Thật sự vậy ạ!?''

Trước khi kịp có người nói thêm gì, tất cả cùng đồng loạt thốt lên trong sự kinh ngạc, giọng nói hòa vào nhau như một bản hợp xướng của sự bất ngờ: ``Hả!?''

\img{7-11p}

Nàng Đeo kính nhìn họ, cũng không kém phần ngơ ngác, đầu cô nghiêng sang một bên với vẻ mặt khó hiểu: ``Sao vậy?''

Sự im lặng lại bao trùm căn phòng vài giây, một khoảng lặng đầy ngượng ngùng. Đây chẳng phải là kết quả mà họ đã mong chờ. Tất cả sự chuẩn bị tinh thần cho một cuộc tranh luận nảy lửa, những kịch bản đối đầu đã được luyện đi luyện lại trong đầu, hóa ra lại vô dụng hoàn toàn, như một vở kịch bị hủy bỏ vào phút chót.

Nàng Thiên sứ cầm tờ giấy báo cáo lên, vẫn chưa thể tin vào tai mình, giọng cô như một tiếng kêu kháng nghị: ``Nhưng ở đây ghi rõ là chúng ta đang có kế hoạch cắt giảm mà!''

Cả lập tức cùng hướng mắt về tờ giấy trên tay cô, như thể đó là một tài liệu mang trọng mạng, một thứ giấy tờ thiêng liêng mà họ đã hiểu sai hoàn toàn.

A-chan nhẹ nhàng nhận lấy tờ giấy, điều chỉnh lại kính rồi đọc kỹ từng dòng chữ, ánh mắt lướt qua những đoạn bị tẩy xóa. Cô im lặng một lát, rồi bỗng ``À!'' lên một tiếng như vừa nhớ ra điều gì, nụ cười ngộ ra hiện rõ trên gương mặt.

Cô vỗ nhẹ vào tờ giấy rồi giải thích với giọng như đang kể một câu chuyện hài hước: ``Các em có biết đây là gì không? Đây là bản kế hoạch ăn kiêng của một nhân viên trong công ty đấy.''

Không khí trong phòng lập tức im lặng đến mức có thể nghe thấy tiếng một cây kim rơi xuống sàn. Một luồng lạnh vô hình thổi qua, và ngay sau đó, tiếng hét vỡ òa bùng nổ.

``\textbf{Hảaaaaa!???}'' cả sáu người đều bật ngửa vì cú sốc, giọng nói như một tiếng sét giữa trời quang. Người thì ôm đầu, người đập bàn, kẻ khác chỉ biết há hốc mồm với vẻ không thể tin vào mắt tai.

Matsuri, mắt đang giật liên hồi, giơ tay lên với giọng như một thám tử đang cố nối các mảnh manh mối: ``Vậy\dots\ con số 100 kia là gì?''

Nàng Đeo kính nhún vai, trả lời một cách tỉnh bơ như đang nói về chuyện hiển nhiên: ``Chắc là số cân mà người đó vượt quá tiêu chuẩn đấy.''

Mio, vẫn còn chưa hoàn hồn, hỏi tiếp với giọng đầy bối rối: ``Còn hạng D thì sao ạ?''

``Chị đoán là kết quả khám sức khỏe tổng quát của người đó.''

Flare đang cố gắng giữ bình tĩnh nhưng mồ hôi đã ướt đẫm trán, giọng cô như một tiếng thở dài não nề: ``Và\dots\ cái thứ cần loại bỏ kia là gì?''

``Mỡ thừa.''

Một tiếng ``phụt!'' vang lên khi Marine suýt nữa phun cả ngụm nước ra, nhưng vẫn cố gắng hỏi thêm với gương mặt chẳng thể tin vào thực tế, giọng cô đầy bất lực: ``Vậy tại sao lại có những đoạn bị tẩy xóa, bị che đi như vậy?''

Nàng Đeo kính bật cười, nụ cười đầy ẩn ý, đôi mắt hiện lên với vẻ thích thú: ``Chị nghĩ chắc người đó thấy ngượng ngùng sau khi viết ra những lời đó thôi.''

\dots Im lặng.

Một khoảng lặng kéo dài vô tận, như thể thời gian đã dừng trôi.

Chỉ còn tiếng tích tắc của đồng hồ và hơi thở phập phồng của sáu con người vừa nhận ra mình đã bị lừa một cách đau đớn. Không có ai bị sa thải. Không hề có cuộc cắt giảm nhân sự nào cả. Tất cả những lo lắng, căng thẳng, những kịch bản đấu tranh mà họ nghĩ ra từ chiều ngày hôm qua, những hoạt động ``tập luyện\dots'' idol mà họ đã bỏ công suy nghĩ thực hiện tới cùng, tất cả đều xuất phát từ\dots\ một bản kế hoạch giảm cân đơn giản bị hiểu lầm.

Và rồi, tiếng hét phẫn nộ bùng lên: ``\textbf{Đúng là lừa gạt mà!}''

\img{7-11q}

Cả sáu người cùng hét lên, mỗi người mang một cảm xúc khác nhau: `` kẻ giận dữ đập bàn, người cười trừ vì thấy bản thân quá ngây thơ, còn kẻ khác chỉ biết thở dài vì đã phí hoài cả buổi sáng - thậm chí là cả ngày hôm qua - vào chuyện vô lý này. Tiếng cười và tiếng than thở hòa quyện, tạo thành một bản giao hưởng của sự bối rối nhưng cũng là nhẹ nhõm.

A-chan chỉ đáp lại bằng một nụ cười chua chát, vừa có chút áy náy, vừa thích thú trước phản ứng của các cô gái, như thể cô đang thưởng thức một vở hài kịch mà mình là khán giả duy nhất. Cô chỉnh kính một lần nữa rồi lẩm bẩm với giọng đầy suy tư: ``Nhưng mà\dots\ tờ giấy này sao lại lạc đến văn phòng của chúng ta nhỉ?''

Ít ai biết rằng, tờ giấy bí ẩn đó vốn bị chủ nhân của nó giấu kỹ trong ngăn bàn, nhưng trong cuộc hỗn loạn ngày hôm qua, nó vô tình bị một cơn gió mạnh thổi từ cửa sổ mở, bay qua mấy dãy bàn làm việc và đáp ngay đống đồ đạc bừa bộn khi mọi người đang dọn dẹp lại phòng ốc - một sự trùng hợp ngẫu nhiên đã tạo nên cả một chuỗi sự kiện hài hước khó quên.

Nhưng ngay sau đó, cô chợt nhớ ra một chuyện quan trọng hơn, giọng cô trở nên hào hứng: ``À, tiện thể các em đang ở đây, bọn chị đang có kế hoạch cho giai đoạn tiếp theo của dự án idol chúng ta này.''

Không khí trong phòng lập tức thay đổi, sự chú ý của cả sáu người được kéo về phía cô như những nam châm. Cả sáu người đều tập trung lắng nghe, ánh mắt đầy háo hức, những gương mặt sáng lên với hy vọng mới.

Nàng Đeo kính tiếp tục, giọng cô đầy nhiệt huyết: ``Ý tưởng lần này là lập các nhóm biểu diễn khác nhau và tổ chức một buổi concert. Chị nghĩ chắc Hikawa-san cũng đã nói qua với các em rồi.''

Vừa nghe đến đó, tất cả họ đều rạo rực, những đôi mắt sáng lên với niềm vui khó tả. Một số gật gù, một số thì lập tức nhớ lại một chuyện quan trọng, như thể một ký ức vừa được đánh thức.

Mio suy nghĩ một chút rồi lên tiếng với giọng như đang nhớ lại một điều gì đó: ``Nhắc mới nhớ, Kuon-senpai cũng đã từng nhắc đến nó rồi\dots''

Matsuri búng tay cái, chợt nhớ ra điều quan trọng với giọng đầy phấn khích: ``À đúng rồi! Lúc mà anh ấy soạn báo giá ấy!''

\begin{center}
\uline{Hồi tưởng\\Một tháng trước.}
\end{center}

Trong văn phòng làm việc, Kuon chăm chú ngồi trước màn hình máy tính đến nỗi quên hẳn mọi chuyện đang diễn ra xung quanh. Trên màn hình là hàng loạt bảng tính chi phí với những con số liên tục thay đổi, phức tạp như một điệu nhảy khó nhằn, bên cạnh còn có một chiếc máy tính cầm tay và nhiều tờ tài liệu đã được ghi chú cẩn thận.

Cậu vừa gõ phím, vừa khẽ lẩm bẩm tính toán, giọng trông giống một nhân viên kế toán đang mệt mỏi vật lộn với các con số: ``\vie{Chi phí dựng sân khấu từ Khoa Kiến trúc công trình\dots\ 100 nghìn yên\dots\ Còn chi phí trang trí như vẽ hoa văn, sơn lại\dots\ tổng cộng khoảng 40 nghìn yên\dots}''

Bất chợt cậu kêu lên một tiếng đầy bất lực: ``\vie{Trời ơi! Tốn nhiều tiền quá vậy!}''

Cậu thở dài, tựa cằm lên tay với vẻ mặt của một người đang phải đối mặt với một ``vực thẳm\dots'' tài chính: ``\textit{Dù huy động cả trường làm mà đã tốn từng này, không biết nếu thuê ngoài sẽ còn mất bao nhiêu nữa\dots\ Với tình hình chi phí đắt đỏ thế này thì phải tiết kiệm thôi\dots}''

Đúng lúc đó, một giọng nói vang lên từ phía sau, làm ngắt quãng dòng suy nghĩ của cậu: ``Anh đang tính toán gì thế?''

``\textbf{Ối giời ơi!!}'' Kuon giật mình đến mức suýt ngã khỏi ghế, tay vội bám chặt vào mép bàn để giữ thăng bằng, rồi quay người lại với vẻ mặt đầy ngạc nhiên.

Cậu thấy Matsuri đứng đó, tay chống hông, mặt tỏ ra tò mò với nụ cười tinh nghịch nở trên môi.

``C-Có chuyện gì à?'' cậu hỏi, vừa cố gắng bình tĩnh lại, tay vô thức đặt lên ngực như để kiểm tra tim đang đập mạnh.

Nàng Cổ động cười lớn, giọng trong trẻo như tiếng chuông lan tỏa khắp phòng: ``Là em đây mà, Matsuri! Em thấy anh ngồi tính toán chăm chú từ nãy đến giờ. Là chuyện gì thế?''

Cậu Tổng nhấc mấy tờ báo giá lên, vẫy vẫy trước mặt với vẻ như đang giải thích một chuyện rất đơn giản: ``À, đây là báo giá để lắp đặt sân khấu và chuẩn bị cơ sở vật chất cho buổi live sắp tới thôi.''

Mio vừa tình cờ đi ngang, nghe vậy liền dừng lại, đôi tai cô lập tức dựng lên vì tò mò: ``Là buổi live vừa rồi của chúng ta hả?''

Cậu Tổng lắc đầu, mắt vẫn dán chặt vào màn hình: ``Không. Là buổi live sắp tới. May thay là tổ chức online nên chi phí rẻ hơn nhiều.''

Nghe vậy, nàng Cổ động mở to mắt, giọng reo lên đầy vui mừng: ``\textbf{Thật vậy á!? Chúng ta sắp có concert hả!?}''

``Mà lại còn là online nữa sao!?'' nàng Sói đen cũng không kém phần ngạc nhiên, đôi mắt sáng bừng lên vì háo hức.

Kuon mỉm cười nhẹ, giọng nói đầy bí ẩn: ``Ừ. Tin chính xác đấy. Ngoài ra còn nhiều điều mới lạ khác nữa. Chính Tanigo-san đã bảo anh soạn cái báo giá này mà.''

Hai cô gái càng lúc càng tỏ ra tò mò, họ nghiêng người về phía cậu như những đứa trẻ đang chờ đợi một câu chuyện cổ tích: ``Cụ thể là thế nào ạ?''

Cậu chỉ lắc đầu với nụ cười đầy ẩn ý: ``Tuyệt mật. Đến ngày hôm đó thì các em sẽ biết thôi.''

Nàng Cổ động bĩu môi, giọng tỏ ra thất vọng: ``Chán thế\~{} Tiết lộ cho người ta một tí mà không được à\~{}''

Dù cậu không tiết lộ nhiều chi tiết, nhưng chỉ cần bấy nhiêu cũng đủ khiến cả hai vô cùng phấn khích, giống như những đứa trẻ đang chờ đợi một món quà bí mật sắp được mở ra.

\begin{center}
\uline{Kết thúc hồi tưởng.}
\end{center}

Sau khi Mio và Matsuri cùng nhau gợi lại kỷ niệm ấy, A-chan gật đầu khẳng định, giọng nói như một người đang dần lắp ráp các mảnh của một dự án lớn thành hình hài hoàn chỉnh: ``Đúng là chị và cậu Hikawa-san đã cùng nhau xây dựng buổi concert này từ đầu, từ lúc lên kế hoạch, họp bàn thảo luận, cho đến khi sân khấu thực sự được dựng nên.''

Rồi cô thốt lên một tiếng cười nhẹ, mang chút tiếc nuối: ``Đáng tiếc là cậu ấy không thể có mặt cùng các em trong buổi biểu diễn này vì có chuyện gia đình. Nhưng chị nghĩ, dù sao thì cậu ấy cũng đã cống hiến rất nhiều cho sự kiện sắp diễn ra này.''

Aqua phấn khích giơ tay lên, giọng reo lên đầy vui mừng: ``Vậy là mọi thứ đã chuẩn bị đâu vào đấy hết rồi sao!?''

A-chan gật đầu, giọng nói đầy tự tin: ``Có thể nói là gần như hoàn thiện cả rồi. Bây giờ chỉ còn bước chia nhóm cho các buổi biểu diễn nữa thôi. Chúng ta sẽ tách ra thành nhiều nhóm khác nhau, lần lượt trình diễn, và tổ chức theo định dạng của một buổi Fes (lễ hội âm nhạc).''

Cả nhóm nghe càng lúc càng hào hứng, đôi mắt ai cũng sáng bừng lên niềm vui chẳng thể nào che giấu. Một buổi concert chia nhóm như thế này chắc chắn sẽ mang lại vô vàn điều thú vị, và chẳng ai trong số họ không muốn tham gia để góp sức của mình.

\img{7-11r}

Matsuri nhiệt tình lên tiếng trước tiên, giọng nói vang dội như một tiếng hô kêu gọi đồng đội: ``Thật ạ? Chị nói thêm chi tiết cho bọn em đi nào!''

Marine cười lớn, tràn đầy háo hức: ``Wow, nghe có vẻ thú vị thật đấy.''

Aqua cũng giơ cả hai tay lên đồng tình, gương mặt rạng rỡ không khác gì một đứa trẻ vừa được tặng kẹo yêu thích: ``Em đồng ý lắm!''

Kanata chợt nảy ra một ý tưởng, mắt cô sáng rực lên: ``Em muốn là chúng ta có thể viết những vở kịch hài nhỏ để trình diễn trong buổi concert ấy!''

Flare lập tức tán thành, giọng nói như một người đang ấp ủ một kế hoạch thú vị: ``Chúng ta có thể lập một nhóm ba người, thế thì sẽ nghĩ ra được nhiều câu đùa vui hơn nữa.''

Mio nhăn mặt, giọng bất lực than thở: ``Lại lại quay sang chuyện hài kịch nữa rồi sao!?''

Mỗi người một ý tưởng, khác nhau hoàn toàn, nhưng may mắn thay lần này A-chan đã ghi lại tất cả một cách cẩn thận, ngọn bút của cô lướt nhanh trên giấy như một vũ công đang biểu diễn trên sân khấu. Cô lắng nghe từng đề xuất, phân tích và đưa ra phản hồi rõ ràng: có cái được chấp nhận ngay, có cái cần chỉnh sửa lại cho phù hợp, cũng có những ý tưởng không thể thực hiện được nên phải loại bỏ.

Tất cả cùng ngồi xuống, những gương mặt rạng ngời háo hức chẳng khác nào những đứa trẻ sắp được nghe một câu chuyện cổ tích kỳ thú. A-chan nhìn các cô với ánh mắt trìu mến, như một người chị lớn đang dõi theo các em nhỏ chuẩn bị bước vào một hành trình mới mẻ.

Cô rút một tập tài liệu từ trên bàn, lật qua những trang giấy chi chít sơ đồ và ghi chú: ``Hiện tại, kế hoạch của chúng ta là chia thành các nhóm nhỏ, mỗi nhóm sẽ đảm nhận một tiết mục khác biệt. Cách tổ chức này sẽ tạo nên sự đa dạng, và cho các em cơ hội thể hiện đúng cá tính của bản thân mình.''

Nàng Hầu gái ngồi sát mép ghế, đôi mắt long lanh đầy phấn khích: ``Vậy là bọn em sẽ được biểu diễn theo nhóm đúng không ạ?''

Nàng Đeo kính gật đầu, giọng đầy chắc chắn: ``Đúng vậy. Các em sẽ được ghép cặp hay lập nhóm dựa trên phong cách riêng của mỗi người. Một số nhóm sẽ mang âm hưởng nhẹ nhàng, lãng đãng, số khác lại có thể sôi động, bùng cháy hơn. Tất cả đều phụ thuộc vào sở trường và mong muốn của từng người.''

Nàng Cổ động nhảy bật lên khỏi ghế, vỗ tay rộn ràng vì vui sướng: ``Ồ! Vậy là bọn em có thể tự chọn người cùng nhóm rồi à?''

Nàng Yêu tinh vàng mỉm cười, đôi mắt ánh lên nét suy tư: ``Có lẽ em và Kanata sẽ cùng làm một tiết mục nhỉ?''

Nàng Thiên sứ nghe vậy, mắt sáng rực lên: ``Ừ! Hai đứa mình sẽ cùng nhau nghĩ ra một tiết mục thật độc đáo, thật đặc biệt!''

Nàng Thuyền trưởng vẫn còn nét tinh nghịch thường ngày, đặt tay lên cằm với gương mặt đầy toan tính: ``Em thì muốn lập nhóm với những người biết cách tạo ra những màn trình diễn\dots\ thật hoành tráng một chút.''

Nàng Sói đen liếc cô bằng ánh mắt ngờ vực, giọng như một lời cảnh báo nhẹ nhàng: ``Hoành tráng kiểu gì đấy? Đừng bảo lại là có màn trượt patin hay xô chậu nhôm như lần trước nữa nhé.''

Cả phòng cùng bật cười, tiếng cười giòn tan lập tức xua tan đi mọi căng thẳng còn sót lại từ ngày hôm qua. A-chan chỉ biết lắc đầu, nhưng khóe môi cũng nở một nụ cười ấm áp.

``Yên tâm đi, bọn chị đã lên kế hoạch kỹ lưỡng hơn thế nhiều,'' cô nói, giọng chứa đựng nét ẩn ý. ``Nhân tiện, nếu các em có ý tưởng gì hay ho, cứ mạnh dạn đề xuất hết. Đây là sân khấu của các em, các em xứng đáng được tỏa sáng theo cách mà các em muốn.''

Matsuri giơ nắm đấm lên cao, giọng đầy nhiệt huyết cháy bỏng: ``Vậy thì chúng ta sẽ tạo ra một buổi concert mà không ai có thể quên!''

Aqua gật đầu mạnh mẽ, đôi mắt tím của cô như bùng cháy ngọn lửa quyết tâm: ``Chúng ta sẽ cùng nhau cố gắng hết sức!''

Kanata cũng hùa theo đầy hứng khởi: ``Tất cả vì khán giả của chúng ta!''

Flare khoanh tay, gật gù đồng tình: ``Có vẻ như chúng ta sắp có thêm một kỷ niệm thật đẹp rồi đây.''

Mio nhìn quanh, bắt gặp những gương mặt rạng rỡ của mọi người, một cảm giác ấm áp tràn đầy trong tim. ``Phải đấy, kỷ niệm đẹp,'' cô lặng lẽ thì thầm với chính mình.

Bầu không khí trong phòng như bừng sáng lên, những tia nắng sớm chiếu qua ô cửa, phủ lên từng gương mặt một lớp ánh sáng ấm áp. Những đề xuất, những kế hoạch và cả những giấc mơ dần dần được hiện thực hóa.

Buổi họp còn kéo dài thêm một lúc. Cùng nhau thảo luận về danh sách bài hát, trang phục biểu diễn, cả những hiệu ứng đặc biệt có thể làm nên điểm nhấn cho tiết mục. Mỗi người đều đóng góp một phần sức mình, dù lớn hay nhỏ.

Và rồi, khi mọi kế hoạch đã được phác thảo xong xuôi, A-chan đứng dậy, khép lại tập tài liệu: ``Vậy là xong. Cảm ơn các em đã dành thời gian hôm nay. Chị tin rằng buổi concert này sẽ là một trong những sự kiện đáng nhớ nhất của \hll.''

Cả nhóm cùng đứng dậy, nụ cười nở trên môi, ánh mắt ai cũng đầy quyết tâm. Họ không còn là những người lo lắng, hoang mang về một tương lai mơ hồ nữa, mà đã trở thành một tập thể vững mạnh, sẵn sàng đón nhận thử thách mới.

Marine bước ra khỏi phòng, ngoái đầu nhìn lại những người còn lại, nụ cười tinh quái nở trên môi: ``Chúng ta sẽ cho họ thấy \hll\ là ai, đúng không nào?''

Tất cả cùng gật đầu, một sự đồng lòng mạnh mẽ như một lời thề hẹn với nhau.

Và thế là, sau những phút giây hỗn loạn, sau bao hiểu lầm và hoảng sợ vô cớ, họ đã tìm thấy điểm chung của mình. Họ sẽ tiếp tục bước đi, cùng nhau, trên con đường mà chính họ đã lựa chọn.

Con đường của những idol không bị gò bó vào khuôn mẫu nào, những idol dám sống đúng với bản thân mình.

%==========================================TRUYỆN PHỤ=========================================
\omk

\pov{Hikawa Kuon}

Mùng 7 Tết Âm lịch.

Đã mười ngày kể từ khi tôi trở về quê ăn Tết. Đây là cái Tết thứ hai kể từ lúc tôi nhận chức Tổng quản lý của \hll.

Nếu tính theo lịch Dương, hôm nay cũng đã là một ngày sau buổi concert online ``Bloom,''. Đáng lẽ đó phải là một buổi diễn tràn ngập tiếng reo hò cổ vũ từ khán giả, nhưng tiếc thay, đại dịch đã khiến họ không thể biểu diễn trực tiếp với khán giả. Những tiếng vỗ tay, những tràng reo hò, và những ánh mắt long lanh của hàng ngàn người hâm mộ - tất cả đều không có mặt trong khán phòng ảo đó.

Chính xác hơn, họ đã phải biểu diễn mà không có khán giả thực sự. Thiếu sự nhiệt tình của người xem, cũng như không có sự tương tác với những người hâm mộ hẳn là một tổn thất lớn đối với các cô gái, những người luôn sống nhờ vào năng lượng từ khán giả. Nhưng với tình hình hiện tại thì chỉ còn cách phải tuân theo thôi, một thực tế mà ai cũng phải chấp nhận.

Dù vậy, dự án này vẫn đánh dấu một bước tiến quan trọng của \hll\ - một dự án còn khá mới mẻ, nhưng đã nhận được sự đầu tư rất lớn cả về chất xám lẫn tài chính. May mắn thay, nhìn vào những phản hồi tích cực từ khán giả, có thể nói nó đã rất thành công, vượt qua cả những kỳ vọng ban đầu của chúng tôi.

Tất nhiên, thành công đó đến từ sự tập luyện chăm chỉ và cống hiến không ngừng nghỉ của các thành viên. Nhưng bên cạnh đó, tôi cũng không thể không cảm ơn những người bạn trong trường đại học của mình. Chính họ đã giúp sức rất nhiều để hoàn thiện phần hậu kỳ cho buổi biểu diễn này, không chỉ nâng cao chất lượng mà còn giúp tôi tiết kiệm được một khoản đáng kể cho công ty.

Thế nhưng, trớ trêu thay, tôi lại chẳng thể theo dõi buổi hòa nhạc ấy một cách trọn vẹn. Tôi không có mặt ở hậu trường dù thật lòng mà nói, tôi cũng không rõ liệu mình có thể tham dự hay không. Ngay cả việc xem trực tiếp qua điện thoại cũng bất khả thi, bởi nhà tôi đang thực hiện chính sách tiết kiệm chi phí một cách khá triệt để - một quy tắc mà tôi chẳng thể nào cãi lại.

Những gì tôi biết về buổi diễn chỉ là qua các tin tức đọc được trên mạng. Tôi có thể hình dung được sự thất vọng của mọi người khi biết tôi đã không theo dõi ``Bloom,'' như họ mong đợi. Dù vậy, khi nghe chị A và đặc biệt là chú Taniho báo tin mừng qua những cuộc gọi ngắn ngủi, tôi cũng đủ hình dung ra nó đã thành công đến mức nào. Tôi gửi lời chúc mừng, cảm ơn mọi người vì sự cống hiến và cũng không quên xin lỗi vì đã bỏ lỡ giây phút quan trọng ấy. Nhưng dù thế nào đi nữa, cuối cùng thì bao nhiêu nỗ lực của chúng tôi cũng đã được đền đáp. Với tôi, chỉ cần vậy là đủ để cảm thấy mãn nguyện rồi.

Họ vui mừng là thế, còn tôi thì sao? Bình thường thôi. Cũng giống như mọi cái Tết mà tôi đã trải qua suốt hơn hai mươi năm nay, một sự lặp lại quen thuộc đến nhàm chán.

Gọi là nghỉ Tết, nhưng thực chất cũng chỉ là danh nghĩa. Từ lúc tôi về đến hết mùng 5, tôi tạm gác công việc sang một bên để tập trung cho gia đình. Nhưng ngay sau đó, tôi lại lao vào kế hoạch tiếp theo cho \hll\ - chủ yếu là trao đổi trực tuyến với chị A và chú Tanigo về những dự án sắp tới, những buổi họp online kéo dài đến khuya. Nếu có ai đó nghĩ tôi là người chăm chỉ\dots\ cũng không hẳn. Vì thật ra, tôi có lý do riêng để làm vậy, một lý do mà có lẽ chỉ tôi mới hiểu.

Nói ra thì có vẻ như đang kể khổ, nhưng với tôi, kể từ khi tôi bước vào thế giới của những công việc, ``ngày nghỉ\dots'' là một khái niệm xa xỉ. Không phải vì công ty bắt ép, mà là do tôi tự muốn như vậy, như một thói quen đã ăn sâu vào máu: `` kiếm thật nhiều tiền về cho gia đình. Đến mức mà ngay cả chú Tanigo - một CEO dày dạn kinh nghiệm - cũng phải lắc đầu vì tôi ``tham việc quá\dots'' dù ông ấy hiểu rằng trong văn hóa làm việc Nhật Bản, chuyện tăng ca chẳng có gì lạ lẫm. Nhưng thói quen của tôi thì có phần quá đà so với tiêu chuẩn thông thường.

Trong suốt một tuần, tôi chỉ nghỉ giỏi lắm là Chủ Nhật, còn các ngày lễ lớn như Tết Dương, thì có thể miễn cưỡng dành ra một ngày. Những dịp mà tôi cho là ``vặt vãnh\dots'' như Tuần lễ Vàng, Ngày Hiến pháp hay Ngày Thiếu nhi, tôi vẫn xung phong làm việc, chẳng bận tâm mọi người nghĩ gì về mình. Có lẽ họ đã quen với sự điên rồ của tôi từ lâu.

Làm việc là để kiếm sống, chẳng phải ai cũng vậy hay sao? Câu hỏi tôi tự đặt ra cho mình mỗi ngày, khi mà cái giá của sự nghỉ ngơi dường như quá đắt so với những gì tôi có thể bỏ lỡ.

Những ngày mà lẽ ra tôi có thể nghỉ ngơi, cộng với gần một tháng nghỉ phép, cuối cùng cũng đủ để tôi an tâm về quê ăn Tết cùng gia đình. Dĩ nhiên, tôi vẫn được trả lương đầy đủ, thế nên ít ra cũng không cần lo lắng về tài chính. Và giờ thì, như các bạn thấy đấy, tôi đang ở nhà đón Tết, trong cái không khí mà tôi đã quen thuộc từ thuở bé.

Nhưng nói thật, Tết năm nay hơi chán. Nếu so với năm ngoái - khi gia đình tôi đón hơn hai mươi người đến chung vui, tạo nên một không khí náo nhiệt không khác gì một hội chợ - thì năm nay rõ ràng ảm đạm hơn hẳn. Nhưng lạ một điều là tôi lại thích sự ``ảm đạm\dots'' này. Nó yên tĩnh, không quá ồn ào, không phải tất bật tiếp khách hay lo toan đủ thứ. Chỉ có điều tôi vẫn cảm thấy thiếu thiếu một chút gì đó, như một nốt nhạc đang bị bỏ lỡ trong một bản giao hưởng.

Không phải vì tôi chán sự lặp lại của ngày Tết, mà chỉ là không có những người đó, không khí có vẻ trống vắng hơn thường lệ. Nhưng mà thôi, tôi còn than thở gì chứ? Chiều tối mai hoặc chậm nhất là sáng ngày kia, tôi đã phải lên đường trở lại thành phố Điện tử để tiếp tục công việc rồi. Vẫn còn một nhiệm vụ quan trọng nữa đang chờ tôi ở phía trước: bàn giao người cuối cùng của nhóm NePoLaBo - Momo\dots\ ờm, gì đó, Nene. Tôi sẽ phải đảm bảo rằng mọi thứ đều suôn sẻ cho cô ấy.

Nếu các bạn thắc mắc tại sao năm nay tôi không mời nhóm \hll\ đến ăn Tết như năm ngoái, thì có ba lý do rõ ràng:

Thứ nhất, gia đình tôi đang thực hiện chính sách tiết kiệm chi phí một cách khá nghiêm ngặt, và mời một đám đông như vậy sẽ là một gánh nặng không nhỏ. Đại dịch cũng khiến việc tụ tập đông người trở nên không an toàn, không phù hợp với các biện pháp phòng chống dịch bệnh mà địa phương tôi đang thực hiện.

Thứ hai, các idol-kiêm-hề cũng đang phải tập luyện gấp rút cho buổi biểu diễn, không có nhiều thời gian để nghỉ ngơi hay đi xa, và tôi không muốn làm gián đoạn lịch trình của họ.

Và quan trọng nhất, nhà tôi năm nay có đám giỗ của ông nội - người mất vào Rằm tháng Giêng năm ngoái\footnote{Xem lại {\flaregothic HISTOIRE 10}, quyển Tết}. Theo phong tục ở đây, trong năm đầu tiên sau khi mất, gia đình sẽ không mời người ngoài đến nhà để tránh mang xui xẻo đến cho họ. Dù mẹ tôi không mấy tin vào những quan niệm xưa cũ này lắm, luôn nói đó chỉ là chuyện cổ lỗ sĩ, nhưng bố tôi lại rất khắt khe, nhấn mạnh rằng phải làm đúng theo lệ tổ tiên để tỏ lòng thành kính, và đó là điều không thể bàn cãi.

Còn một lý do nữa, tuy không phải là chính, nhưng chắc chắn có sức nặng. Đó là Mano Aloe đang sống chung với chúng tôi (các bạn đã đọc ở quyển D thì chắc các bạn sẽ hiểu). Nếu các thành viên \hll\ mà kéo đến thì ``kế hoạch của chúng ta coi như hỏng bét!'', hoặc ít nhất đó sẽ là lời của chú Tanigo. Thực ra tôi cũng không rõ cái kế hoạch đó cụ thể là gì, nhưng nếu nó đổ bể thì chắc chắn không phải là chuyện hay ho gì cả, và tôi không muốn là người phải gánh chịu hậu quả.

Nhắc đến cô nàng Succubus ấy, đây là cái Tết đầu tiên em ấy được đón cùng gia đình tôi sau khi chuyển đến ở chung. Mấy ngày vừa qua, em ấy cũng hòa mình hoàn toàn vào không khí ngày Tết, từ cùng mẹ tôi đi chợ mua sắm, phụ làm các món ăn truyền thống, đến tham gia tất cả các nghi lễ của gia đình. Không khí nhà vì có thêm em ấy mà càng thêm ấm áp, náo nhiệt, như thể từ trước đến nay em ấy vốn là một thành viên không thể thiếu của gia đình.

Ban đầu tôi dự định sẽ nghỉ đến hết Rằm để dự đám giỗ gia đình, nhưng có lẽ tôi sẽ quay trở lại thành phố sớm hơn để tiếp tục công việc. Một guồng quay mới lại bắt đầu, và dự cảm của tôi cho thấy năm nay tôi sẽ còn bận rộn hơn cả hai năm trước cộng lại.

Nhưng đó là chuyện của tương lai, còn bây giờ thì sao? Bây giờ, tôi chỉ muốn tận hưởng những phút giây yên bình cuối cùng trước khi lại lao vào vòng xoáy công việc.

\il

Quay trở về thực tại.

Tôi cũng đang viết báo cáo cho công ty, giống như mọi hôm. Ngoại trừ không hoàn toàn như vậy, vì lần này tôi đang phải cố hoàn thành nó trong khi xung quanh chẳng có lấy một chút không khí làm việc. Mọi người trong nhà vẫn còn tận hưởng những ngày Tết cuối cùng, tiếng cười nói và tiếng chén đũa va vào nhau vang lên từ phòng bên, còn tôi thì dán mắt vào màn hình, cố gắng hoàn thành nốt công việc trước khi trở lại thành phố Điện tử. Những âm thanh ấm áp của gia đình như một bức tường ngăn cách tôi với thế giới bên ngoài, nhưng tôi biết mình sẽ sớm phải rời đi.

Trong lúc tôi đang cặm cụi gõ phím với những con số và biểu đồ, bỗng màn hình máy tính báo có cuộc gọi video đến. Nhìn vào màn hình, tôi nhận ra ngay người gọi: `` giám đốc Tanigo, với khuôn mặt quen thuộc hiện lên trong khung hình nhỏ.

``\vie{Chắc là gọi hỏi về công việc đây mà,}'' tôi lẩm bẩm, khẽ thở dài rồi bấm nhận cuộc gọi, một cảm giác vừa quen thuộc vừa mệt mỏi tràn qua.

``À, cậu Hikawa-san đấy hả?'' Đúng, là ông ấy rồi, giọng nói trầm ấm đặc trưng vang lên từ loa.

Tôi vẫy tay chào lại, cố gắng giữ vẻ tự nhiên: ``Vâng, chào chú ạ.''

``Mọi việc thế nào rồi?'' Đúng, tôi nghĩ là ông ấy đang hỏi tôi về công việc đây, như thường lệ.

``À, cái bản báo cáo cháu soạn sắp xong rồi, tầm tối thì cháu--\dots'' tôi nói, nhưng chưa kịp dứt lời thì ông ấy đã cắt ngang. ``Không, không, tôi không hỏi công việc. Tôi hỏi việc nhà cậu ấy.''

Tôi hơi bất ngờ, suýt nữa thì lỡ nhịp đánh máy. Bình thường, mỗi khi chú ấy gọi gọi thì hầu hết là để bàn về công việc, chứ hiếm khi lại hỏi han chuyện cá nhân như thế này. Một sự thay đổi bất thường khiến tôi cảnh giác.

Dù tôi cảm thấy hơi nghi ngờ, nhưng tôi vẫn nhún vai trả lời với giọng bình thản: ``Nhà cháu ạ? Nhà cháu bình thường. Mọi người đều khỏe.''

``Ăn Tết vui vẻ chứ?''

``\dots Vâng\dots? Chú cũng có thể nói như vậy.'' Tôi trả lời một cách chung chung, vừa nói vừa nhìn ra cửa sổ, nơi những cành đào vẫn còn rực rỡ sắc hồng. Đúng là Tết năm nay có hơi tĩnh lặng hơn so với năm ngoái, nhưng nếu bảo tôi đánh giá nó có vui hay không thì cũng khó nói. Có lẽ tôi đã quá quen với sự bận rộn để cảm nhận được niềm vui trọn vẹn.

``Cậu nói thế là tôi thấy yên tâm rồi.''

``Thế ạ? Vậy để cháu cúp máy--\dots'' tôi di chuột về nút kết thúc cuộc gọi với một cảm giác nhẹ nhõm, nhưng ông ấy chợt thốt lên: ``Ấy ấy, từ từ, cái cậu này! Đã nói xong đâu!''

``\dots Dạ?'' tôi nghiêng đầu khó hiểu, tay dừng lại giữa chừng.

``Tôi gọi cho cậu đâu chỉ mỗi hỏi thăm sức khỏe đâu. Họ cũng muốn nói chuyện với cậu đấy.''

``Họ? Ý chú là\dots'' Một cảm giác bất an bắt đầu len lỏi trong lòng tôi.

Nghe có gì đó không ổn lắm. Có lẽ ông ấy đang không nói về những người đã cố gắng gọi cho tôi lúc tôi đang nghỉ Tết đấy chứ? Tôi đã cố tình tránh những cuộc gọi đó, và bây giờ họ đã tìm cách tiếp cận tôi qua một con đường khác.

``Cậu đoán đúng rồi đấy. Đã lâu ngày rồi cậu chưa về, các cô gái giờ cũng nhớ cậu lắm.''

Tôi cau mày, một chút hoài nghi hiện trên gương mặt: ``Chú nói thế nào chứ trông họ vẫn vui mà. Ở trên live show họ diễn hôm qua ấy, cháu biết thừa.''

``Bên ngoài khác với bên trong đấy, cậu bé à. Thôi, để tôi dẫn họ nói chuyện với cậu.''

Nói rồi, ông giám đốc rời khỏi chỗ ngồi, màn hình rung nhẹ khi ông ấy di chuyển. Tôi không biết mình có nên vui hay lo lắng nữa, một cảm xúc hỗn độn đang dâng lên trong lòng.

Tôi nghe thấy tiếng ông ấy gọi mấy người vào, những giọng nói quen thuộc vang lên từ xa. Trong lúc đó, tôi tranh thủ rời khỏi chỗ ngồi để lấy cốc nước, vừa để tránh bị gọi ngay lập tức, vừa để có thời gian chuẩn bị tinh thần. Những bước chân của tôi vang lên nhẹ nhàng trên sàn nhà, như một lần trì hoãn vô ích.

Mà thực ra, cũng có gì đâu mà phải chuẩn bị?

Tôi không thể gọi đó là cảm giác khó chịu, nhưng thú thực là tôi vẫn chưa quen với những cuộc trò chuyện thế này. Một phần vì tính tôi không thích chen chúc trong những nơi đông người - live concert dù là online cũng không ngoại lệ. Đó là lý do tôi quyết định nghỉ Tết ở nhà, thay vì ở lại công ty để tham dự hay theo dõi trực tiếp các buổi diễn của họ. Không phải tôi không quan tâm, chỉ là tôi không giỏi hòa vào bầu không khí ấy, như một cây xương rồng lạc giữa rừng hoa.

Thật ra, nếu là những nơi như xe buýt hay tàu điện - những không gian chật hẹp nhưng ai cũng có mục đích riêng, ai làm việc nấy - thì tôi lại chẳng thấy phiền. Nhưng concert thì khác. Ở đó, mọi người hòa mình vào dòng cảm xúc chung, cùng reo hò, cổ vũ, chia sẻ niềm vui với nhau. Còn tôi? Tôi chỉ cảm thấy như một người ngoài cuộc, một khán giả vô hình đứng nhìn từ một góc xa xôi. Tôi đã thử vài lần, nhưng cuối cùng vẫn không thể thay đổi suy nghĩ đó, vẫn là cái cảm giác lạc lõng giữa đám đông.

Nhấp một ngụm nước, tôi quay trở lại máy tính. Và ngay lập tức, đập vào mắt tôi là cảnh sáu người cùng hô lên với một nụ cười rạng rỡ, những khuôn mặt tươi tắn như những bông hoa mùa xuân: ``\textbf{Kuon-senpai! Chúc mừng năm mới!}''

``Ối chà, mấy đứa quý hóa ghê nhỉ?'' tôi đáp lại, đặt cốc nước xuống với một nụ cười gượng gạo, cố tỏ ra thoải mái dù trong lòng đang cảm thấy hơi choáng ngợp.

Mio hỏi thăm với giọng đầy quan tâm: ``Mấy ngày nay anh ăn Tết với gia đình thế nào?''

``Ờm\dots\ cũng vui.'' Có lẽ tôi không biết phải nhận xét Tết năm nay như thế nào thật, chỉ đưa ra một câu trả lời khá hờ hững.

``Mọi người ở nhà anh vẫn khỏe chứ?'' Matsuri hỏi, giọng đầy háo hức. Này này, mấy đứa định hỏi lại đúng những câu mà Tanigo-san vừa hỏi anh đấy à? Tôi không khỏi cảm thấy một chút bối rối trước sự đồng điệu của họ.

Nhưng dù vậy, tôi vẫn giữ một nụ cười, điềm nhiên đáp với giọng bình thản: ``Ổn. Rất ổn. Cảm ơn.''

Tôi và các cô gái tiếp tục câu chuyện trong cuộc gọi video, tiếng cười và những câu chuyện vui vẻ vẫn không ngừng vang lên. Đến lượt nàng Yêu tinh vàng, nhưng không phải là hỏi thăm sức khỏe của tôi: ``Em đã được các senpai ở trên kể về trải nghiệm ăn Tết quê anh rồi. Nghe có vẻ vui lắm nhỉ?''

``Anh thấy nó cũng na ná giống Tết Dương lịch ở bên đó thôi. Nhưng, ừm, nó vui.'' Tôi nhấp một ngụm nước, cố gắng tìm cách diễn tả một cách đơn giản nhất.

``Nó như thế nào vậy?'' Kanata chen vào, tỏ vẻ tò mò, đôi mắt cô sáng lên như một đứa trẻ đang chờ nghe một câu chuyện cổ tích.

``Nói bây giờ ra hơi khó giải thích, nhưng thôi được, để anh tóm gọn lại nhé\dots''

Thế là tôi bắt đầu kể lể về những nét đặc trưng của Tết quê tôi, từ những món ăn truyền thống không thể thiếu đến các nghi thức độc đáo của ngày lễ. Mọi người đều chăm chú lắng nghe, đặc biệt là Marine, Flare và Kanata, những cái đầu cứ liên tục gật gù theo từng câu chuyện tôi kể. Những cô gái khác dù đã từng trải qua qua Tết cùng tôi rồi, nhưng vẫn không giấu được sự hào hứng, cứ như đang nghe một câu chuyện cổ tích thú vị nào đó lần đầu tiên.

``Hôm nay em mới biết ở quê anh còn có tục lệ thờ cúng tổ tiên nữa nha,'' nàng Thiên sứ ngây ngô nói, vẻ mặt như vừa khám phá ra một điều gì đó thật thú vị.

``Đó là chuyện bình thường mà,'' tôi nhún vai, trả lời một cách tự nhiên.

``Thế có giống như là cầu phúc cho những người đã mất lên thiên đường không ạ?''

``Cũng\dots\ tương tự vậy thôi.'' Tôi chẳng hiểu sao em ấy lại dùng cách ví von đó, nhưng cũng không muốn giải thích dài dòng.

``Em cũng mới biết đến Ông Công Ông Táo đó,'' Matsuri chen vào, đôi mắt long lanh đầy ngạc nhiên, chẳng khác nào một đứa trẻ vừa tìm ra một điều mới lạ.

Tôi bật cười trước phản ứng đáng yêu của cô: ``Nhân vật này nổi tiếng đến mức có cả chương trình truyền hình lấy làm chủ đề luôn. Ví dụ như chương trình `Gặp nhau cuối năm' ấy, chuyên kể lại những sự kiện đáng chú ý trong một năm. Chỉ có điều là phải cập nhật tin tức thường xuyên mới hiểu hết được những câu chuyện trong đó. Chương trình đó còn có cả những bài hát nhép rất vui nhộn nữa.''

Vừa khi tôi dứt lời, Aqua lập tức hào hứng quay sang nhìn tôi, đôi mắt sáng rực, giọng cô trong trẻo như tiếng chuông reo: ``A, nhắc đến nhạc thì\dots\ Anh đã xem buổi biểu diễn của chúng em chưa ạ!?''

\dots Chết thật. Tôi đứng hình ngay lập tức, chẳng biết nên trả lời sao nữa. Biết rồi mà sẽ có người hỏi câu này, như một quả bom hẹn giờ vừa kích nổ đúng lúc.

Nói thật là tôi sợ làm các em buồn, mà nếu dối trá thì lại sợ bị các em tra khảo thêm. Cứ thế này thì chắc chắn họ đã nghi ngờ tôi không xem buổi diễn rồi. Chết tiệt, thực sự là khó để thoát khỏi tình thế này, tôi cứ như một con thú đang bị mắc vào bẫy vậy.

May mắn thay, đúng lúc đó, chị đồng nghiệp của tôi bước ra từ phía sau nhóm và nói giúp tôi, giọng cô đầy sự bao dung: ``Hikawa-san bận rộn với việc gia đình nên cậu ấy không thể theo dõi được.''

Cảm ơn chị thật nhiều, A-san! Chị vừa cứu tôi khỏi một tình thế thiệt thòi! Tôi thở phào nhẹ nhõm trong lòng, tựa như một người vừa thoát khỏi cơn bão tố đột ngột.

Đúng như tôi dự đoán, tất cả các cô gái đều tỏ ra hơi buồn, nhưng họ cũng rất thông cảm với hoàn cảnh của tôi. Những gương mặt có chút buồn thảm, nhưng rồi lại nhanh chóng nở nụ cười khi thấy vẻ áy náy hiện rõ trên mặt tôi.

Chị ấy nói cũng có phần đúng. Mùng 1 Tết cha, mùng 2 Tết mẹ, chưa kể còn đám giỗ của ông nội tôi sắp tới nữa. Tôi chạy cả ngày ngoài đường để thăm họ hàng cũng đủ thấm mệt rồi, lấy đâu ra sức để mà xem nọ xem kia chứ. Những ngày Tết của tôi là những ngày của những cuộc gặp gỡ và những chuyến đi dài.

Chị A nhìn tôi và nói, giọng cô đầy dịu dàng: ``Chúng tôi cũng đã ghi lại toàn bộ buổi concert `Bloom,' ấy rồi, khi nào cậu đi làm trở lại sẽ đưa cho cậu, coi đó như là quà Tết của công ty cho cậu ha.''

\dots Chắc là cũng nên nhận thôi nhỉ? Dù chẳng biết tôi có bao giờ xem chúng không. Tôi gật đầu với vẻ biết ơn: ``Dạ, vâng, em xin nhận.''

``Thế là may rồi nhỉ, Kuon-senpai?'' Aqua nhìn tôi với sự nhẹ nhõm trên mặt, như thể vừa tránh được một thảm họa. ``Giờ anh không lo sẽ bị bỏ lại phía sau nha!''

``Bọn em có nhiều trò vui trong buổi concert ấy lắm! Khi nào anh về chúng ta cùng xem lại nha!'' Matsuri nói với tông giọng vui vẻ, giọng cô như một lời hứa hẹn về những giây phút vui vẻ sắp tới.

Tôi khẽ gật đầu, mỉm cười: ``Được chứ, anh chờ mấy đứa đấy.''

Và rồi, mười phút tiếp theo trôi qua trong niềm vui chung của tất cả mọi người, chúng tôi cùng ngồi lại chia sẻ những câu chuyện hài hước mà cả nhóm đã từng trải qua. Tiếng cười vang lên rộn rã qua loa máy tính, như thể các cô gái đang ngồi ngay bên cạnh tôi chứ không phải ở đầu dây bên kia.

Một trong những câu chuyện đó lại là một sự nhầm lẫn đáng nhớ mà chắc chắn sẽ đi vào lịch sử của công ty. ``Thế hóa ra mấy đứa lại đọc nhầm cái gì vậy? Tờ tuyên ngôn ăn kiêng mà hiểu lầm thành thông báo sa thải á?''

Mio lập tức phồng má lên, phụng phịu như một đứa trẻ vừa bị trêu chọc: ``Cũng đâu có vui đâu chứ! Tờ giấy đó bị gạch xóa lung tung chi chít, mà chữ viết của người đó thì\dots\ khó đọc kinh khủng luôn!''

``Ồ, thế sao? Đưa anh xem nào?''

Ngay lập tức, Matsuri và Marine cùng giơ lên một tờ giấy đầy những nét gạch ngang, gạch chéo đan xen. Dòng chữ ban đầu viết khá nắn nót, nhưng lại bị tô đè lên bằng những nét bút nguệch ngoạc. May mắn là camera của họ khá rõ nét, tôi vẫn có thể đọc được từng chữ, nhìn những đường gạch xóa như thể thấy được cuộc đấu tranh nội tâm của người viết.

``À, anh biết cái này của ai rồi! Đúng là `tuyên ngôn ăn kiêng' của Usake-san!'' Chỉ cần nhìn lối viết đặc trưng đó là tôi nhận ra ngay, thứ mà tôi đã quá quen thuộc.

``Thế rốt cuộc chuyện là thế nào vậy ạ?'' Flare nghiêng đầu hỏi, ánh mắt đầy tò mò nhìn vào tờ giấy trên màn hình.

``Ừm\dots\ cách đây không lâu, ông ấy than thở với anh rằng dạo này trông mình hơi mập ra. Anh hỏi cân nặng thì ông ấy bảo: `Tôi có 90 cân thôi mà.' Thật không đùa đâu, bản thân anh cũng nặng tới 95 cân, nhưng anh cao hơn, còn ông ấy thì thấp hơn anh một chút. Chính sự khác biệt đó mà câu chuyện mới trở nên buồn cười.''

Nhớ lại cảnh đó, tôi vừa cười vừa kể tiếp: ``Anh liền bảo: `Thế là vừa đẹp mà, cần gì ăn kiêng?' Nhưng ông ấy lại lắc đầu quầy quậy: `Không được! Kỳ này tôi phải ăn kiêng cho bằng được!' Và thế là\dots\ tờ tuyên ngôn này ra đời.''

``Rồi cuối cùng kết quả sao rồi anh?'' Marine bụm miệng cười hỏi, giọng trông háo hức chờ đợi phần kết của câu chuyện hài.

``Nhịn được đúng ba ngày. Đến ngày thứ tư thì kêu đói quá không chịu được, thế là\dots\ thôi, bỏ cuộc luôn!''

Cả sáu cô gái cùng lúc ngửa người ra sau, bật cười nghiêng ngả, tiếng cười vọng qua loa như một bản nhạc vui nhộn. Tôi cũng nói thêm: ``Anh nghĩ ông ấy gạch xóa chi chít như thế là sợ bị phát hiện, mà cũng đâu cần, muốn giấu thì bỏ vào máy hủy giấy là xong mà!''

Aqua cười đến mức nước mắt rơi, giọng cô nghẹn lại vì quá vui: ``Trời ơi, thế mà cũng dám tuyên bố ăn kiêng nữa!''

Matsuri vừa cười vừa đập bàn, tiếng đập bàn vang lên như nhịp trống: ``Không có chí tiến thủ gì hết trơn á!''

``Thì nói gì đến Mat-chan, năm ngoái em cũng có lần thử nhịn ăn ăn kiêng y hệt thế\footnote{Xem lại {\flaregothic ÉPISODE 33}, quyển số Ba.} mà suýt ngất xỉu đó. Nhớ hồi đó không?''

Bị gọi tên, nàng Cổ động chỉ biết cười trừ, má cô hơi đỏ lên vì xấu hổ: ``T-Tại hồi đó em chưa biết ăn kiêng đúng cách mà\dots''

Tôi lại hỏi tiếp, tỏ vẻ hiếu kỳ với tờ giấy ấy: ``Mà này, sao mấy đứa lại có được tờ giấy này? Cũng tình cờ thôi à? À mà hồi anh bật máy tính lên hôm qua thì thấy Flare-chan có gọi điện và nhắn tin cho anh đó.''

``Đúng vậy đó! Nhưng mà Mio-senpai bảo anh không bao giờ nhấc máy, cũng chẳng chịu trả lời tin nhắn. Sao lại thế ạ?'' Flare hỏi, giọng có chút trách móc nhẹ nhàng, dù môi cô vẫn đang nở nụ cười.

``À, đúng là anh đã thế! Đã nghỉ Tết thì phải để công việc sang một bên chứ! Anh chỉ giữ liên lạc với A-san và Tanigo-san thôi. Mà với hai người đó, anh cũng dặn trước rồi: chỉ khi nào họp toàn công ty hay có việc cực kỳ khẩn cấp mới gọi, bình thường thì cứ để anh yên.'' Đó là quy tắc mà tôi đặt ra từ đầu năm nay, sau một vài sự kiện bất ngờ, để mình tôi có được chút không gian riêng tư.

``Nhưng nếu là việc của bọn em thì sao? Em còn tưởng là anh bỏ rơi bọn em rồi đấy\dots'' nàng Thiên sứ bất mãn nói, giọng như một đứa trẻ đang trách móc người lớn bỏ mặc mình.

Tôi chỉ nhún vai một cách thản nhiên: ``Thì phải tự giải quyết với nhau chứ. Anh đâu thể lúc nào cũng ở bên các em để xử lý hết mọi chuyện đâu? Mấy đứa cũng phải lớn lên, phải trưởng thành thôi.''

``Nhưng mà sau đó bọn họ đều bắt nạt em hết!'' Kanata rít lên, giọng đầy ấm ức như đang kêu cứu. Nhìn vẻ mặt là biết cô ấy đang chất chứa nhiều chuyện không vừa lòng.

``Sao thế? Kể anh nghe xem có chuyện gì.''

Ngay lập tức, nàng Thiên sứ phàn nàn đủ thứ, kể lại việc bị lôi ra làm ``chuột bạch\dots'' cho đủ thứ ý tưởng tập luyện thần tượng kỳ quái. Từ việc phải nhảy dây liên tục trong khi hát, đến tập múa với tạ đeo ở chân. Thậm chí cái vầng hào quang đặc trưng của em ấy còn bị lấy ra làm đồ chơi, ném qua ném lại như một chiếc đĩa bay.

Tôi bật cười lớn đến mức không thể ngăn lại: ``Trời ơi, tưởng có chuyện gì kinh khủng lắm mới gọi cứu hộ! Mấy đứa mà chả hay bày đủ trò quái đản này suốt ngày, đúng là đám hề chính hiệu của công ty, không ai đọ được!''

Chẳng đợi kịp ai cất lời, tiếng cười của cả nhóm lại ùa tới, vang vọng qua loa máy tính như một bản hòa tấu vui nhộn. Matsuri cười đến mức phải ôm bụng, còn Marine thì vừa lau nước mắt cười vừa trách: ``Chuyện đó đâu phải lỗi tại bọn em! Ai để Kanata cứ tỏ ra thiệt thòi thế chứ, ai thấy cũng phải xúi giục cho hết!'' Nàng Thiên sứ nghe vậy chỉ biết bĩu môi trêu lại, làm cả nhóm lại cười lớn thêm một trận nữa. Những tiếng đập bàn xen lẫn tiếng cười nói làm tan biến đi chút ngượng ngùng vừa có, mọi người dường như quên hết mọi lo âu để tận hưởng khoảnh khắc thoải mái này.

Sau khi cười đùa đủ, cả nhóm mới từ từ bình tĩnh lại, bầu không khí ấm áp vẫn còn đó nhưng đã nhẹ nhàng hơn hẳn. Tôi vẫn để nụ cười trên môi, quan sát từng gương mặt qua màn hình: Kanata vẫn còn hơi phụng phịu, còn Aqua thì đang khoanh tay với vẻ đầy tự hào, như thể em ấy vừa giấu được một điều lớn lao nhất định phải chia sẻ với tôi.

``Em đã định nghĩa với mọi người thế nào mới là một idol đúng nghĩa rồi đấy anh!'' nàng Hầu gái nói, giọng đầy hãnh diện như vừa khám phá ra một chân lý vĩ đại.

Tôi nhướng mày, tò mò hỏi: ``Ồ? Vậy định nghĩa của em là gì nào?''

Aqua nhấp một ngụm nước, hắng giọng một cách chỉnh tề như một diễn giả chuẩn bị lên sân khấu, rồi bắt đầu nói với giọng đầy khí thế: ``Một idol mà em mơ ước trở thành, đó là người được làm những gì mình đam mê, và quan trọng hơn hết, có thể cống hiến hết mình để mang niềm vui đến cho mọi khán giả. Đó là quan điểm của em, và ai cũng có quyền xây dựng con đường idol của riêng mình theo cách họ muốn!''

Giọng nói của cô nàng chắc chắn, không hề có chút do dự nào. Tôi khẽ gật đầu, cảm nhận được sự chân thành và nghiêm túc tỏa ra từ lời nói của em ấy, một điều mà tôi hiếm khi thấy ở Aqua thường ngày hay đùa giỡn.

``Đúng rồi, chuẩn không cần chỉnh. Làm idol đâu nhất thiết phải là người đứng trên sân khấu lớn, phải tỏa sáng dưới hàng ngàn ánh đèn sân khấu hào nhoáng mới đúng? Aqua-chan nói quá đúng rồi.''

Matsuri khoanh tay, nhìn tôi với ánh mắt nửa đùa nửa thật: ``Ê, hôm nay lạ nhỉ, đến cả Kuon-senpai cũng đồng ý với ý kiến của em ấy. Hiếm có thật sự!''

Tôi chỉ cười khẽ. Cái kiểu nói chuyện cà khịa, trêu chọc này của mấy đứa thì đúng là năm tháng trôi qua cũng chẳng bao giờ thay đổi. Nhưng lần này tôi không để em ấy trêu chọc mãi, mà tiếp tục lái câu chuyện vào chủ đề chính, giọng nói của tôi cũng dần trở nên nghiêm túc hơn.

``Nhưng anh phải nhấn mạnh thêm, chắc mấy đứa cũng hiểu rồi, rằng dù con đường trở thành idol của mỗi người có khác nhau thế nào đi nữa, thì trách nhiệm đi cùng với danh hiệu đó luôn rất nặng nề, không thể xem thường.''

``Điều đó bọn em đều hiểu rõ,'' Mio lên tiếng trước, giọng cô chắc nịch như đại diện cho cả nhóm khẳng định, ánh mắt sắc bén toát lên sự quyết tâm không hề đùa giỡn.

Tôi nhìn lướt qua từng người, thấy ai cũng đã nghiêm túc lắng nghe, bèn tiếp tục chia sẻ: ``Theo anh nghĩ, idol chính là thần tượng, là một hình mẫu mà người khác muốn noi theo, là một con người gương mẫu mà mọi người có thể học hỏi điều tốt đẹp. Nhưng idol không nhất thiết phải là người nổi tiếng trên mạng, hay phải đứng trên sân khấu lớn. Idol, hay thần tượng, có thể đơn giản là bố mẹ trong gia đình, một người bạn thân, hay thậm chí là thầy cô giáo đã dìu dắt mình.''

Tôi ngừng lời một lát, để những câu nói đó có thời gian thấm vào tâm trí của từng cô gái. Có những điều nghe thì đơn giản, nhưng phải có người nhắc lại thì mới thực sự ngẫm được hết ý nghĩa, như một bài học cuộc sống mà chỉ khi được nhắc nhở mới thật sự khắc sâu vào lòng.

Marine là người đầu tiên lên tiếng, giọng cô có chút băn khoăn: ``Vậy\dots\ ý của anh là\dots?''

``Ý anh là, anh không bao giờ áp đặt các em phải xây dựng hình tượng idol theo một khuôn khổ nào. Đó là quyền lựa chọn của riêng mỗi người. Nhưng điều anh luôn mong mỏi, là các em sẽ xây dựng con đường đó sao cho vừa phù hợp với cá tính của chính mình, đồng thời luôn giữ được sự văn minh, tinh thần trách nhiệm, và không bao giờ đi ngược lại những quy tắc đạo đức tốt đẹp của xã hội. Các em hiểu chứ?''

Tất cả mọi người đồng thanh gật đầu, những cái gật đầu đầy quyết tâm. Không cần nói thêm lời nào, tôi biết các em đã hoàn toàn hiểu được thông điệp mà tôi muốn gửi gắm.

Tôi dựa lưng chậm rãi vào ghế, nhìn vào màn hình nơi những gương mặt quen thuộc của các thành viên hololive vẫn đang hiện lên. ``Idol\dots'' không chỉ đơn thuần là một nghề, hay một cái danh hiệu để khoe, mà nó còn là một lý tưởng, một trách nhiệm lớn lao mà các em đã tình nguyện gánh vác trên vai. Và tôi thật sự hy vọng, những lời chia sẻ này sẽ giúp các em vững bước hơn trên con đường mình đã chọn.

``Làm idol là một hành trình rộng lớn, đầy tự do và niềm vui\dots\ nhưng đi cùng với đó cũng là vô vàn áp lực. Đợi đến khi các em trưởng thành thêm một chút nữa, các em sẽ thấu hiểu hết những điều anh nói ngày hôm nay.''

Đó chính xác là những bài học về con đường idol mà tôi muốn gửi gắm cho tất cả các em, nhân dịp năm mới Tân Sửu này. Tôi nhìn những gương mặt đang chăm chú lắng nghe, và tôi tin chắc rằng, dù phía trước có bao nhiêu thử thách chờ đợi, các em sẽ luôn vững bước, vì các em đã có những người đồng hành tuyệt vời bên cạnh để cùng vượt qua mọi khó khăn.

Chúng tôi còn rất nhiều chuyện để chia sẻ, nhưng đáng lẽ ra cuộc gọi hôm nay sẽ có thể kéo dài\dots

\dots nếu không có \emph{một người nào đó} bất ngờ xuất hiện và phá tan mọi thứ.

``\textbf{Kuon-kun\~{}!! Bánh quy em làm xong rồi nè! Anh muốn thử một cái không!?}'' một giọng nói trong trẻo, đầy hứng khởi vang lên từ cầu thang, như một tiếng chuông reo làm rung chuyển không khí.

Tôi đứng hình hoàn toàn, toàn thân cứng ngắc như pho tượng đá. Trời ơi, là Aloe!

Cô nàng Succubus hiện ra khỏi bóng tối, tay bưng một khay bánh quy vừa ra lò, mặt rạng rỡ tự hào về thành quả của mình. Những chiếc bánh vàng ươm xếp gọn gàng trên khay, chỉ chờ được thưởng thức. Nhưng ngay khi giọng nói của em ấy vang lên, bầu không khí trong cuộc gọi lập tức thay đổi. Từ màn hình máy tính, tôi có thể thấy những ánh mắt nghi ngờ lần lượt bật sáng, như những ngọn nến lần lượt được thắp lên trong đêm tối.

Aloe nghiêng đầu, nhìn vào màn hình với vẻ tò mò, đôi mắt to tròn ngây thơ hỏi: ``Anh đang nói chuyện với ai đấy, Kuon-kun? Để em xem với nào--\dots''

Tôi vội vàng nhảy ra chắn trước em ấy, mồ hôi toát ra như mưa, tim đập thình thịch như trống làng đang hội: ``K-Không có gì đâu!''

Ở phía bên kia màn hình, Kanata thắc mắc nhướng mày, giọng đầy tò mò: ``Hả? Ai vậy, Kuon-senpai?''

Aqua thì nheo mắt lại đầy nghi ngờ, giọng của cô nàng như một thám tử đang thẩm vấn nghi phạm: ``Anh có bạn gái rồi à?''

Tôi ấp úng, cổ họng khô rang như thể vừa nuốt hết cát sa mạc: ``Ờm\dots\ không đâu\dots\ chỉ là bạn bè thôi\dots''

Đương nhiên lời giải thích yếu kém đó chẳng thể nào thuyết phục được họ. Marine càng lúc càng hào hứng trêu chọc tôi, giọng cô đầy tinh nghịch kéo dài: ``Anh có bạn gái rồi ư\~{}?''

Tôi đứng bật dậy, vừa cố gắng che chắn Aloe sau lưng vừa phản bác lại mạnh mẽ: ``Không có đâu! Mấy đứa nghĩ gì lung tung thế!''

Ngay cả Aqua và Matsuri cũng không đứng về phía tôi. Cả hai cùng đồng thanh nhìn tôi, rồi lại nhìn bóng dáng người phía sau với ánh mắt đầy nghi vấn: ``Trông có vẻ không đơn giản đâu nhỉ\dots''

Nhìn vẻ mặt tò mò sục sạo của bọn họ, tôi biết nếu tiếp tục cuộc gọi này, tôi chắc chắn sẽ bị tra hỏi cho đến khi lộ hết mọi chuyện. Sự thật sẽ chẳng thể nào giấu được nữa.

``Thôi được rồi, mấy hôm sau anh đi làm lại ta kể tiếp nhé! Tạm biệt các em!'' Tôi lập tức tắt cuộc gọi trước khi ai kịp phản ứng gì, ngón tay bấm nút kết thúc một cách vội vàng như đang né tránh một quả bom sắp nổ.

Trời ơi, Aloe ơi là Aloe\dots\ em ấy vẫn chưa sẵn sàng để gặp mặt những người kia đâu\dots

Tôi chớp mắt, cảm nhận không khí xung quanh bỗng im lặng đến lạ, chỉ còn tiếng gió thổi lướt qua khe cửa. Chắc hẳn sáu người ở đầu dây bên kia đang vô cùng bối rối. Mà không chỉ họ, ngay cả bản thân tôi cũng cảm thấy hành động của mình hơi đột ngột quá. Tại sao lại phải cúp máy vội vàng như vậy chứ?

Vài giây trôi qua, tôi thở dài một hơi, quay sang nhìn Aloe. Cô nàng vẫn đứng đó, chớp mắt ngây thơ như chẳng hiểu chuyện gì, đôi mắt long lanh lộ vẻ bối rối. Tôi lại thở dài một lần nữa, giọng đầy bất lực: ``\dots Lần sau em đừng có hét toáng lên như thế nữa được không?''

Aloe nhíu mày, nghiêng đầu trông vô tội hết mực, giọng như một đứa trẻ đang biện minh cho hành động của mình: ``Nhưng em đã làm bánh quy cho anh mà! Em muốn anh được nếm thử thành quả em tự tay làm mà!''

Ngay lúc đó, một ánh mắt sắc bén khác cũng hướng vào tôi. Ryuuhou đã đứng khoanh tay quan sát từ nãy giờ, cuối cùng cũng mở lời với giọng đầy nghi vấn: ``\dots Anh vừa gọi điện thoại với ai thế, Onii-san?''

Tôi im lặng trong giây lát, lướt mắt sang nhìn Aloe. Tôi không thể nói ra sự thật khi em ấy vẫn còn ở đây, bí mật về công việc của tôi phải được giữ kín.

Từ góc độ của cô nàng Succubus ấy, chắc chắn em ấy nghĩ tôi đang che giấu điều gì đó mà em không được phép biết, một sự bí ẩn đáng ngờ.

Nhưng thay vì để bản thân bị kéo vào cuộc tra hỏi không hồi kết, tôi thản nhiên lấy một chiếc bánh quy, cắn một miếng, rồi gật đầu với vẻ mặt như thể chẳng có chuyện gì xảy ra: ``Không có gì đâu, ăn bánh thôi.''

Chưa kịp để Aloe phản ứng gì, Ryuuhou cũng với lấy một chiếc bánh, nhấm nháp thử với vẻ mặt dè dặt. ``Ừm\dots\ cũng ngon đấy.''

Đôi mắt của cô nàng lập tức sáng rực như đèn pha vừa được bật, nghe lời khen bất ngờ đó em ấy reo lên: ``Thật hả!?''

Và chỉ với một câu nhận xét ngắn gọn đó, mọi chuyện vừa xảy ra lập tức bị quẳng vào quên lãng. Aloe nhanh chóng ngồi xuống bên chúng tôi, hào hứng chia sẻ về món bánh mà em đã tự tay làm, tiếng cười của cô vang vọng khắp phòng như một bản nhạc vui nhộn.

Tôi lén nhìn em ấy một chút, sau đó lặng lẽ nhâm nhi bánh quy cùng với tách trà nóng, cố gắng tỏ ra như thể chẳng có gì đặc biệt vừa xảy ra. Hy vọng là sáu người ở đầu dây bên kia không để ý quá nhiều đến vị khách bất ngờ kia.

Cầu trời cầu phật cho điều đó.

\il

\dots Được rồi, tôi phải thừa nhận rằng tình huống vừa rồi khá là hỗn loạn. Suýt chút nữa là mọi bí mật của tôi với chú Tanigo đã bị bại lộ hoàn toàn.

Mà khoan, lúc nãy tôi định nói gì để kết thúc câu chuyện này nhỉ? À đúng rồi, e hèm.

Mùa đông cũng dần đi qua. Một mùa xuân mới đang đến gần, mang theo không khí tươi mới và vô vàn điều thú vị đang chờ đợi. Điều đó cũng có nghĩa là, núi công việc lại đang chờ tôi tại văn phòng của \hll\ này.

Sau buổi concert ``Bloom,'' sẽ còn rất nhiều buổi biểu diễn khác nữa. Sẽ có thêm nhiều thành viên mới gia nhập đại gia đình, sẽ có thêm những câu chuyện vui, cũng như những khoảnh khắc đáng nhớ, cả buồn lẫn vui.

Nhưng dù có thế nào đi nữa, tôi vẫn sẽ là một phần của gia đình \hll\ - một nơi đầy rẫy những tình huống éo le, những chuyện điên rồ, nhưng cũng tràn đầy sự ấm áp và tiếng cười.

Và tôi sẽ tiếp tục cố gắng hết sức mình!

Chúc mừng năm mới Tân Sửu!

\pov{-}

Ở phía đầu dây bên kia, ngay sau khi màn hình cuộc gọi tắt đi một cách đột ngột, cả sáu cô gái vẫn đứng hình, những gương mặt còn đang đóng băng trong sự bất ngờ.

Matsuri chớp mắt mấy lần để tỉnh táo, rồi búng tay phát ra tiếng ``tách\dots'' đầy hứng khởi, ánh mắt cô rạng rỡ như một thám tử vừa tìm ra mảnh ghép quan trọng của vụ án: ``Thế mới hay chứ! Chẳng phải Kuon-senpai vừa bị người ta bắt quả tang còn gì nữa?''

``Ừ thì đúng, nhưng mà là ai thế nhỉ?'' Aqua nghiêng đầu, tay chống cằm ra chiều suy luận, bộ óc cô dường như đang quay cuồng với vô số giả thuyết. ''Giọng nói nghe thì\dots\ trẻ trung, nhưng mà cảm giác hơi quen quen sao á?''

``Đã nói là bánh quy mà!'' Mio bất ngờ cất tiếng, khiến tất cả mọi người lập tức quay đầu nhìn về phía cô. Cô gật gù như vừa chộp được một mảnh quan trọng của câu đố, đôi mắt sắc bén như lưỡi liềm: ``Cô bé kia hét lên gì mà 'Bánh quy em làm xong rồi đây!!' đúng không? Nghe y hệt vậy.''

``Chậc, vậy là người quen của Kuon-senpai à?'' Marine vừa nói vừa huých nhẹ vào vai Flare, nhưng cô nàng Yêu tinh chỉ nhún vai với vẻ mặt bối rối, chẳng thể nêu ra được cái tên nào: ``Không nhớ ra là ai rồi, nhưng mà giọng đó chắc chắn mình đã nghe qua đâu đó\dots''

Kanata chống nạnh, đôi cánh nhỏ sau lưng rung rung theo nhịp suy nghĩ, giọng cô đều đều như đang đọc lại các đặc điểm của một kẻ tình nghi trong danh sách: ``Một cô gái có giọng nói trẻ con, năng động, biết làm bánh quy\dots\ hừm\dots''

Tất cả sáu người nhìn nhau, cùng nhau vắt óc lục lọi trong kho trí nhớ để mò ra cái tên đó. Cái tên\dots\ có chữ ``Ma\dots'' gì đó à? Hay là bắt đầu bằng chữ ``A\dots''? Nó cứ lơ lửng đâu đó trong đầu, như một bong bóng xà phòng, muốn chạm tới lại vụt mất.

``\dots Mà này khoan, lúc nãy anh ấy hoảng loạn thật sự đúng không?'' nàng Cổ động đột nhiên reo lên, đôi mắt long lanh như vừa khám phá ra một bí mật động trời. ``Chị thề là từ trước đến nay chưa bao giờ thấy Kuon-senpai phản ứng kiểu đó đâu! Chắc chắn là có chuyện mờ ám ẩn giấu!''

Nàng Hầu gái liên tục gật đầu như cái máy, giọng cô đầy thuyết phục như thể vừa nhìn thấy bằng chứng thép: ``Chắc chắn 100\%! Biểu cảm lúc đó của anh ấy là kiểu 'Chết thật, lộ hết rồi' á, nhìn rõ rành rành không thể nhầm lẫn!''

Trong khi nhóm cô gái vẫn đang bàn tán sôi nổi, xô xát nhau để đưa ra giả thuyết, thì ở một góc phòng khác, A-chan thở dài một hơi, tay đẩy gọng kính lên với vẻ mặt bất lực. Cô quay sang nhìn người đàn ông đứng cạnh mình, người cũng đang nở một nụ cười khó xử không kém.

``Tanigo-san\dots\ kế hoạch của ông suýt nữa thì bị bại lộ hoàn toàn rồi đấy.''

Ông chủ của công ty chỉ cười trừ, khoanh tay nhìn vào màn hình cuộc gọi giờ đã trống trơn, giọng nói đầy vẻ bí ẩn: Nhưng cuối cùng thì đâu có bị lộ ra đâu, đúng chứ?''

Nàng Đeo kính lắc đầu bất lực, nhưng cũng chẳng nói thêm câu gì nữa. Dù sao thì, đây cũng chẳng phải lần đầu tiên có chuyện như vậy xảy ra, và chắc chắn cũng sẽ không phải là lần cuối cùng.

Cuối cùng, chàng quản lý đã phải đứng ra giải thích hết cho mọi chuyện, nói dối rằng đó là một cô em họ của cậu sang chơi, tới lúc đó mọi người mới thôi không tra hỏi thêm nữa, dù rằng họ vẫn còn rất nhiều câu hỏi trong đầu, những nghi vấn chưa có lời giải đáp.

Không biết liệu năm Tân Sửu này Hikawa Kuon còn có những tam tai nào nữa đây? Chỉ có thời gian mới trả lời được.

\begin{center}{\ab BONUS SONG}\\(Đây đơn thuần chỉ là phần hát của các thành viên sẽ xuất hiện trong buổi live Bloom, đó. Bạn có thể xem video ở dưới, nếu không thích đọc phần tự sự của tôi.)
  \vid{7-11s}{song}\\(Đặt giả thuyết là tất cả mọi người đều cùng hát bài hát đó.)
\end{center}

Câu hát mở đầu với hai câu:

\begin{center}
  少しずつ近づいているよ　夢に見たトキメキ \\
  (Niềm háo hức mà chúng ta hằng mong chờ đang đến gần)

  広がってく世界中に \\
  (Tất cả, chúng được rải ra trên khắp thế gian này)
\end{center}

Cùng với đó, tất cả các thành viên \hll\ đều dần dần xuất hiện, bắt đầu di chuyển (dạng hình tĩnh) từ phải sang: Roboco, Mel, Aki, Matsuri. Câu hát tiếp tục:

\begin{center}
  皆を笑顔にさせるキラメキを \\
  (Chúng ta đây sẽ tỏa sáng lấp lánh, khiến cho mọi người đều mỉm cười)

  一人じゃないからライバルがいるから \\
  (Vì chúng ta không hề đơn độc, vì chúng ta đều có đối thủ của nhau)
\end{center}

Tiếp tục là Fubuki, Aqua, Shion và Mio.

\begin{center}
  競い合うことすら \\
  (Dù cho là một cuộc chiến khốc liệt đến cỡ nào)

  楽しくなっちゃうんだよ \\
  (Nó cũng sẽ trở thành một niềm vui)
\end{center}

Kế đến là Ayame, Subaru, Choco và Miko.

\begin{center}
  {\fotskip Dear My Friends} \\
  (Hỡi người bạn của tôi)

  未来へ走り出したら{\fotskip Fantastic} \\
  (Khi ta hướng tới tương lai, mọi thứ sẽ rất là tốt đẹp ở phía trước)

  夢を咲かせて{\fotskip Our Future} \\
  (Ta sẽ cùng biến những ước mong nở rộ của mình vào tương lai của chúng ta)

  一歩一歩を踏み出して　トップを目指そう \\
  (Từng bước, từng bước, ta từ từ hướng tới đỉnh cao)

  信じる力が心の{\fotskip Happiness} \\
  (Hãy cùng tin tưởng vào bản thân, để rồi sẽ đưa con tim mình trở nên hạnh phúc)
\end{center}

Tiếp theo là năm thành viên thế hệ thứ Ba: Pekora, Rushia, Marine, Flare và Noel.

\begin{center}
  全力で楽しもう \\
  (Ta cùng tận hưởng những gì mà ta có)

  それはいつかきっと \\
  (Bởi chắc chắn một ngày nào đó)
\end{center}

Cuối cùng là bốn thành viên mới nhất từ \hll\ thế hệ thứ năm: Botan, Polka, Lamy và Nene\footnote{Sẽ được giới thiệu ở quyển số Tám.}

\begin{center}
  大きな夢叶えるから \\
  (Chúng ta sẽ biến những ước mơ to lớn đó trở thành sự thật)

  {\fotskip Dreaming Days} \\
  (Hỡi những ngày mơ mộng của chúng ta\dots)
\end{center}

Đoạn phần ca hát đó kết thúc, màn hình quảng cáo cho Bloom, xuất hiện, cùng với đó là câu kết:

\begin{center}
  \textbf{Đây sẽ là nơi mà chúng ta nở rộ, ngay tại chính nơi này\dots}
\end{center}

\end{document}